%%%%%%%%%%%%%%%%%%%%%%%%%%%%%%%%%%%%%%%%%%%%%%%%%%%%%%%%%%%%%%%%%%%%%
%% Supplementary Material -- Spectrochimica Acta Part A
%% Standalone document. Compile with pdfLaTeX. No bibliography.
%% Sections, tables, figures and equations are all S-prefixed.
%%%%%%%%%%%%%%%%%%%%%%%%%%%%%%%%%%%%%%%%%%%%%%%%%%%%%%%%%%%%%%%%%%%%%
\documentclass[11pt]{article}

\usepackage[T1]{fontenc}
\usepackage[margin=1in]{geometry}
\usepackage{amssymb}
\usepackage{amsmath}
\usepackage{float}
\usepackage{graphicx}
\usepackage{xcolor}
\usepackage{booktabs}
\usepackage{array}
\usepackage{longtable}
\usepackage{url}

\newcommand*\wn{\ensuremath{\text{cm}^{-1}}}
\providecommand{\SI}[2]{#1\,#2}

\renewcommand{\thesection}{S\arabic{section}}
\renewcommand{\thesubsection}{S\arabic{section}.\arabic{subsection}}
\renewcommand{\thetable}{S\arabic{table}}
\renewcommand{\thefigure}{S\arabic{figure}}
\renewcommand{\theequation}{S\arabic{equation}}
\setcounter{secnumdepth}{2}
\setcounter{tocdepth}{2}
\makeatletter
\renewcommand{\@seccntformat}[1]{\csname the#1\endcsname\quad}
\renewcommand{\@pnumwidth}{2.4em}
\renewcommand{\@tocrmarg}{3.0em}
%% S-prefixed numbers are wider than the article defaults allow
\renewcommand*{\l@section}[2]{%
  \ifnum \c@tocdepth >\z@
    \addpenalty\@secpenalty
    \addvspace{1.0em \@plus\p@}%
    \setlength\@tempdima{2.9em}%
    \begingroup
      \parindent \z@ \rightskip \@pnumwidth
      \parfillskip -\@pnumwidth
      \leavevmode \bfseries
      \advance\leftskip\@tempdima
      \hskip -\leftskip
      #1\nobreak\hfil
      \nobreak\hb@xt@\@pnumwidth{\hss #2}\par
    \endgroup
  \fi}
\renewcommand*{\l@subsection}{\@dottedtocline{2}{2.9em}{3.3em}}
\makeatother

\title{Supplementary Material\\[0.4em]
       \large The computed library, not the spectroscopy, limits Raman identification of authentic chemical warfare agents}
\author{%
\begin{minipage}{0.95\textwidth}\centering
Doo-Hee Lee\textsuperscript{a,b,$\dagger$},
Seongchan Lee\textsuperscript{a,$\dagger$},
Sungbong Kim\textsuperscript{c},
Jinkwang Jeong\textsuperscript{c},\\
Rosy Oh\textsuperscript{d},
Hyein Kim\textsuperscript{a},
Ku Kang\textsuperscript{a,*}\\[0.9em]
\normalsize
\textsuperscript{a}Chemical, Biological, Radiological and Nuclear (CBRN) Defense Research
Institute, Ministry of National Defense, Seoul 06796, Republic of Korea\\[0.15em]
\textsuperscript{b}Department of Civil Engineering and Environmental Sciences,
Korea Military Academy, Seoul 01805, Republic of Korea\\[0.15em]
\textsuperscript{c}Department of Chemistry, Korea Military Academy,
Seoul 01805, Republic of Korea\\[0.15em]
\textsuperscript{d}Department of Mathematics, Korea Military Academy,
Seoul 01805, Republic of Korea\\[0.7em]
\footnotesize
\textsuperscript{$\dagger$}Doo-Hee Lee and Seongchan Lee contributed equally to this work.\\
\textsuperscript{*}Corresponding author: bisu9082@gmail.com
\end{minipage}}
\date{}

\begin{document}


\maketitle

\tableofcontents
\clearpage

%%%%%%%%%%%%%%%%%%%%%%%%%%%%%%%%%%%%%%%%%%%%%%%%%%%%%%%%%%%%%%%%%%%%%
\section{Corrections to the previous analysis}
\label{si:corrections}

This work has been analyzed three times. Each reanalysis changed results that had
already been reported, and we gather every correction here rather than distributing
them through the text.

\subsection{Corrections made in the second analysis}

\textbf{The calibration was selected on the metric it reported.} In the first analysis
the frequency scaling factor and the linewidth were chosen by a grid search whose
objective was the number of agents identified correctly at rank~1, which was the quantity
then reported as the result. Every calibration since obtains the scale factors from band
positions and the linewidth from measured band widths, neither of which is the reported
metric.

\textbf{The band-position statistics were inflated by the matching scheme.}
Nearest-neighbor matching allowed one computed band to be claimed by several measured
bands, with residuals of opposite sign that partly canceled.

\textbf{A coefficient of determination was reported for the parity comparison.} It is
uninformative over a 1550\,\wn{} window, since uniform noise at the matching tolerance
reproduces it to three decimals, and it has been removed instead of recomputed.

\textbf{Three substantive claims were withdrawn.} That fitting the linewidth to spectral
similarity degrades identification; that the measured A-232 spectrum contains no resolved
bands below 700\,\wn{}; and that the intense 700--760\,\wn{} feature is the P--F stretch,
which normal-mode analysis shows to be a mixed motion of the F--P--C fragment with the
largest P--F character lying higher.

\textbf{The A-series structures were corrected before the present calculations.} An
earlier geometry set drew the amidine terminal amine as
\emph{N}-methyl-\emph{N}-ethyl rather than \emph{N,N}-diethyl. All four A-series
calculations were repeated from corrected structures.

\subsection{Corrections made in the present analysis}

The second analysis worked from spectra that had already been background-corrected and
rescaled by a procedure we could characterize but not specify. A referee identified that
dependency as disqualifying. The present analysis begins at the instrument's raw detector
counts, and rebuilding the pipeline from there changed five things that had been assumed.

\textbf{Twenty-three of the fifty-five delivered files are not acquisitions.} They
reproduce a genuine file to a multiplicative factor near unity with an added offset and
noise, and none carries integer detector counts. Every count in the second analysis
treated them as independent measurements. HD and VX have a single genuine acquisition
each, not five (Table~\ref{tab:si-sessions}).

\textbf{Integration time is recorded per file and varies by a factor of 252.} It was
assumed to be unrecorded and common across agents. It is neither, and it is a confounder
between agents.

\textbf{The exclusion of HN2 was justified on the wrong grounds.} The second analysis
attributed it to a poor measurement. HN2 has the highest signal-to-noise ratio of the
eleven agents and the second highest count rate (Table~\ref{tab:si-quality}). The
exclusion stands, but on spectral content rather than acquisition quality, and
Section~\ref{si:ident} gives independent evidence that the material is not HN2.

\textbf{The measurement grid is 270 points, not 1024.} The detector carries 512 pixels,
of which 270 fall in the 250--1800\,\wn{} fingerprint region. The 2.27\,\wn{} spacing
assumed previously is a product of vendor interpolation. A band-position residual of
10.4\,\wn{} is approximately 1.8 pixels, not 6.
   
\textbf{Three detector pixels are defective and two lie in diagnostic windows}
(Section~\ref{si:preproc}). The region-resolved failure analysis of the second analysis
was computed with them included.

\textbf{Two claims of the second analysis are withdrawn.} That removing the
600--850\,\wn{} window accounts for every molecular error: masking the window rescues one
agent and sends two others across class boundaries. And that the deficit in that window
is attributable to conformational sampling: the most flexible agent is the best in it and
the most rigid the worst (Section~\ref{si:ident}).

\clearpage

%%%%%%%%%%%%%%%%%%%%%%%%%%%%%%%%%%%%%%%%%%%%%%%%%%%%%%%%%%%%%%%%%%%%%
\section{Raw acquisitions and the preprocessing specification}
\label{si:preproc}

\subsection{What the delivered files contain}

Table~\ref{tab:si-sessions} is the inventory. A file is counted as a genuine acquisition
if its RawData column contains integer detector counts, which is what the instrument
writes; the derived files fail that test and are excluded throughout. Nine agents retain
three to five genuine acquisitions and two retain one.

%%%%%%%%%%%%%%%%%%%%%%%%%%%%%%%%%%%%%%%%%%%%%%%%%%%%%%%%%%%%%%%%%%%%%
%%  table: sessions
%%%%%%%%%%%%%%%%%%%%%%%%%%%%%%%%%%%%%%%%%%%%%%%%%%%%%%%%%%%%%%%%%%%%%
\begin{table}[htbp]
\centering
\caption{Inventory of raw acquisitions. Genuine acquisitions are those whose RawData column contains integer detector counts; the derived files reproduce a genuine parent to a multiplicative factor near unity with an added offset and noise, and are excluded throughout. Integration time is recorded per file by the instrument. Bands are those detected above a normalized intensity of 0.10 in the 250--1800\,\wn{} region.}
\label{tab:si-sessions}
\footnotesize
\resizebox{\textwidth}{!}{%
\begin{tabular}{lrrrrrrr}
\toprule
Agent & Genuine & Derived & Integration (ms) & Peak counts & SNR & Counts\,s$^{-1}$ & Bands \\
\midrule
GD & 4 & 1 & 1161--1936 & 34783 & 420 & 27194 & 20 \\
GF & 3 & 2 & 834 & 35376 & 491 & 42418 & 22 \\
VX & 1 & 4 & 1700 & 32511 & 317 & 19124 & 21 \\
HD & 1 & 4 & 25995 & 10668 & 36 & 410 & 7 \\
HN1 & 3 & 2 & 1014 & 36217 & 454 & 35717 & 17 \\
HN2 & 3 & 2 & 119 & 34193 & 643 & 287336 & 8 \\
HN3 & 3 & 2 & 5840 & 34517 & 88 & 5910 & 7 \\
A-230 & 3 & 2 & 20851--30000 & 10590 & 28 & 392 & 10 \\
A-232 & 3 & 2 & 5445 & 35369 & 302 & 6496 & 18 \\
A-234 & 5 & 0 & 1471--4937 & 29004 & 345 & 18611 & 19 \\
A-242 & 3 & 2 & 5356--24634 & 23176 & 153 & 3674 & 15 \\
\bottomrule
\end{tabular}
}

\begin{flushleft}\footnotesize Total: 32 genuine acquisitions of 11 agents from 55 files. HD and VX have a single genuine acquisition each and are never used as queries.\end{flushleft}
\end{table}
%%%%%%%%%%%%%%%%%%%%%%%%%%%%%%%%%%%%%%%%%%%%%%%%%%%%%%%%%%%%%%%%%%%%%
\clearpage

\subsection{Data-quality criterion and the exclusion of HN2}

A criterion was fixed before any identification was attempted: a measured spectrum is
retained if at least four bands reach 10\% of its strongest band within
250--1800\,\wn{}. Table~\ref{tab:si-quality} applies it.

%%%%%%%%%%%%%%%%%%%%%%%%%%%%%%%%%%%%%%%%%%%%%%%%%%%%%%%%%%%%%%%%%%%%%
%%  table: quality
%%%%%%%%%%%%%%%%%%%%%%%%%%%%%%%%%%%%%%%%%%%%%%%%%%%%%%%%%%%%%%%%%%%%%
\begin{table}[htbp]
\centering
\caption{Data-quality criterion, applied before any identification was attempted. A measured spectrum is retained if at least four bands reach 10\% of its strongest band within 250--1800\,\wn{}.}
\label{tab:si-quality}
\small
\begin{tabular}{lrrrrrl}
\toprule
Agent & Bands & $\geq$10\% & $\geq$20\% & SNR & Top-band share & Decision \\
\midrule
GD & 20 & 15 & 7 & 420 & 0.190 & retained \\
GF & 22 & 15 & 9 & 491 & 0.210 & retained \\
VX & 21 & 20 & 17 & 317 & 0.100 & retained \\
HD & 7 & 7 & 7 & 36 & 0.280 & retained \\
HN1 & 17 & 8 & 4 & 454 & 0.260 & retained \\
HN2 & 8 & 1 & 1 & 643 & 0.810 & \textbf{excluded} \\
HN3 & 7 & 7 & 6 & 88 & 0.260 & retained \\
A-230 & 10 & 10 & 10 & 28 & 0.180 & retained \\
A-232 & 18 & 16 & 8 & 302 & 0.170 & retained \\
A-234 & 19 & 13 & 10 & 345 & 0.170 & retained \\
A-242 & 15 & 13 & 11 & 153 & 0.180 & retained \\
\bottomrule
\end{tabular}

\begin{flushleft}\footnotesize HN2 fails the criterion, and the reason is spectral content rather than acquisition quality: its signal-to-noise ratio is the highest of the eleven agents and its count rate the second highest. The previous version of this work attributed the exclusion to a poor measurement, which these figures do not support. Section~S6 gives the independent evidence from the computed library.\end{flushleft}
\end{table}
%%%%%%%%%%%%%%%%%%%%%%%%%%%%%%%%%%%%%%%%%%%%%%%%%%%%%%%%%%%%%%%%%%%%%
\clearpage

\subsection{The specification}

Each genuine acquisition was processed as follows. Dark counts were subtracted using the
instrument's own dark column. Cosmic rays were removed by replicate consensus, described
below. Three defective detector pixels were masked and interpolated. The spectrum was
trimmed to 250--1800\,\wn{} on its native axis, with no interpolation onto a finer grid.
The background was removed by asymmetric least squares with $\lambda = 10^4$ and
$p = 0.01$, values fixed from the literature before any identification was run. Intensities
were normalized to the maximum, with standard-normal-variate normalization reported
alongside throughout.

\textbf{Cosmic rays are removed by replicate consensus.} Flagging narrow positive
excursions one spectrum at a time marks 7.8 pixels per acquisition, 1.5\% of the
512-pixel detector. Of those flags 81.1\% reproduce at the same pixel in an independent
acquisition of the same agent. A cosmic ray is a random event and cannot reproduce; what
reproduces is signal. Interpolating only pixels flagged in one acquisition and not in its
replicates reduces the removal to 1.8 pixels per acquisition, 0.34\% of the detector. The
conventional single-spectrum treatment would have deleted six pixels of real signal per
acquisition. HD and VX have one genuine acquisition each and cannot be treated this way;
for those two the residual contamination is bounded above by the 0.34\% rate measured on
the others.

\textbf{Three detector pixels are defective.} A criterion fixed before the identification
analysis identifies pixels 136, 206 and 286, at 813.7, 1220.6 and 1634.4\,\wn{}. A
pixel qualifies if its deviation from a local baseline is negative in at least 90\% of
acquisitions and exceeds six local standard deviations in at least half of them. The first lies at
the upper edge of the 600--850\,\wn{} window, the second inside the 1150--1400\,\wn{}
window and the third at the foot of the C=N stretching region. A sign fixed across every
acquisition and every agent is an instrumental property. Masking and interpolating the
three raises the separation between within-class and between-class similarity from 0.090
to 0.122.

\textbf{The remaining analyst freedom is bounded.} Table~\ref{tab:si-preproc} gives the
full grid.

%%%%%%%%%%%%%%%%%%%%%%%%%%%%%%%%%%%%%%%%%%%%%%%%%%%%%%%%%%%%%%%%%%%%%
%%  table: preproc
%%%%%%%%%%%%%%%%%%%%%%%%%%%%%%%%%%%%%%%%%%%%%%%%%%%%%%%%%%%%%%%%%%%%%
\begin{table}[htbp]
\centering
\caption{Sensitivity of the preprocessing specification. Cosine is taken against the adopted setting (bold); separation is the mean within-class minus the mean between-class cosine over the eleven measured agents. Entries for the fingerprint limits change the grid itself, so no cosine is defined.}
\label{tab:si-preproc}
\small
\begin{tabular}{llrrrr}
\toprule
Parameter & Value & Cosine vs adopted & Within & Between & Separation \\
\midrule
AsLS lambda & 1000.0 & 0.9960 & 0.607 & 0.485 & 0.122 \\
AsLS lambda & 3000.0 & 0.9987 & 0.614 & 0.489 & 0.125 \\
AsLS lambda & \textbf{10000.0} & 1.0000 & 0.619 & 0.496 & \textbf{0.123} \\
AsLS lambda & 30000.0 & 0.9973 & 0.625 & 0.506 & 0.119 \\
AsLS lambda & 100000.0 & 0.9863 & 0.636 & 0.522 & 0.113 \\
AsLS lambda & 1000000.0 & 0.9666 & 0.642 & 0.539 & 0.104 \\
AsLS p & 0.001 & 0.9821 & 0.659 & 0.559 & 0.100 \\
AsLS p & 0.005 & 0.9984 & 0.631 & 0.516 & 0.115 \\
AsLS p & \textbf{0.01} & 1.0000 & 0.619 & 0.496 & \textbf{0.123} \\
AsLS p & 0.02 & 0.9988 & 0.605 & 0.474 & 0.131 \\
AsLS p & 0.05 & 0.9924 & 0.579 & 0.440 & 0.139 \\
fingerprint low & 200.0 & n/a & 0.616 & 0.496 & 0.120 \\
fingerprint low & 250.0 & n/a & 0.619 & 0.496 & 0.123 \\
fingerprint low & 300.0 & n/a & 0.625 & 0.497 & 0.128 \\
fingerprint low & 400.0 & n/a & 0.628 & 0.499 & 0.129 \\
fingerprint high & 1600.0 & n/a & 0.620 & 0.495 & 0.124 \\
fingerprint high & 1700.0 & n/a & 0.619 & 0.496 & 0.124 \\
fingerprint high & 1800.0 & n/a & 0.619 & 0.496 & 0.123 \\
normalization & \textbf{max} & 1.0000 & 0.619 & 0.496 & \textbf{0.123} \\
normalization & snv & 0.8254 & 0.495 & 0.306 & 0.188 \\
\bottomrule
\end{tabular}

\begin{flushleft}\footnotesize The cosine against the adopted spectrum is not defined for the rows that change the analysis window, because a different window puts the spectrum on a different grid; those cells read n/a. The baseline parameters were fixed at values conventional in the literature before any identification was run. Over three orders of magnitude in $\lambda$ the spectra stay within a cosine of 0.967 of the adopted setting and the separation moves between 0.104 and 0.125; over $p$ from 0.001 to 0.05 the cosine stays above 0.982 and the separation moves between 0.100 and 0.139. Standard-normal-variate normalization gives a higher separation (0.188) while changing the spectra substantially, and both are reported rather than the better-scoring one being adopted.\end{flushleft}
\end{table}
%%%%%%%%%%%%%%%%%%%%%%%%%%%%%%%%%%%%%%%%%%%%%%%%%%%%%%%%%%%%%%%%%%%%%
\clearpage

%%%%%%%%%%%%%%%%%%%%%%%%%%%%%%%%%%%%%%%%%%%%%%%%%%%%%%%%%%%%%%%%%%%%%
\section{Computational library}
\label{si:library}

\subsection{Composition}

The library contains 54 computed Raman spectra. Eleven are the measured agents, including
HN2, which failed the data-quality criterion, and VX, which yields a single genuine
acquisition; both are retained in the library without being queried. The remainder are
organophosphorus analogs of the general form
R\textsuperscript{1}O--P(R\textsuperscript{2})(=O)F. GF (cyclosarin) is simultaneously a
measured agent and the cyclohexyl/methyl member of that analog grid, and is counted
once. Forty-six of the 54 members are nerve-class, 4 are vesicants and 4 are Novichok
agents; class assignments in the main text are made against that composition.
Table~\ref{tab:si-library} lists every member with the SMILES string from which its
starting geometry was built.

The SMILES strings of the measured agents were checked against the elemental composition
of the corresponding optimized geometry; all agree.

%%%%%%%%%%%%%%%%%%%%%%%%%%%%%%%%%%%%%%%%%%%%%%%%%%%%%%%%%%%%%%%%%%%%%
%%  table: library
%%%%%%%%%%%%%%%%%%%%%%%%%%%%%%%%%%%%%%%%%%%%%%%%%%%%%%%%%%%%%%%%%%%%%
{\small
\begin{longtable}{ll l >{\ttfamily\footnotesize}p{5.4cm}}
\caption{Composition of the 54-member computed spectral library. GF occupies the cyclohexyl/methyl cell of the analog grid and is therefore both a measured agent and a grid member; it is counted once. HN2 was measured but failed the data-quality criterion (Table~\ref{tab:si-quality}) and was retained in the library without being queried; VX has a single genuine acquisition and is likewise retained without being queried.}\label{tab:si-library}\\
\toprule
ORCA job & Molecule & Role & SMILES \\
\midrule
\endfirsthead
\multicolumn{4}{l}{\footnotesize\emph{Table \ref{tab:si-library}, continued}}\\
\toprule
ORCA job & Molecule & Role & SMILES \\
\midrule
\endhead
\bottomrule
\endfoot
EXP\_A230 & A-230 & measured agent & {CP(=O)(F)N=C(C)N(CC)CC} \\
EXP\_A232\_Eglob & A-232 & measured agent (global-minimum geometry, S9) & {COP(=O)(F)N=C(C)N(CC)CC} \\
EXP\_A234 & A-234 & measured agent & {CCOP(=O)(F)N=C(C)N(CC)CC} \\
EXP\_A242 & A-242 & measured agent & {CP(=O)(F)N=C(N(CC)CC)N(CC)CC} \\
EXP\_GD\_soman & GD (soman) & measured agent & {CC(C)(C)C(C)OP(C)(=O)F} \\
EXP\_HD & HD (sulfur mustard) & measured agent & {ClCCSCCCl} \\
EXP\_HN1 & HN1 & measured agent & {CCN(CCCl)CCCl} \\
EXP\_HN2 & HN2 & measured agent & {CN(CCCl)CCCl} \\
EXP\_HN3 & HN3 & measured agent & {ClCCN(CCCl)CCCl} \\
GF\_cyclosarin\_glob & GF (cyclosarin) & measured agent (global-minimum geometry, S9) & {CP(=O)(F)OC1CCCCC1} \\
EXP\_VX & VX & measured agent (not queried; one acquisition) & {CCOP(C)(=O)SCCN(C(C)C)C(C)C} \\
EtGB & EtO--P(Me)(=O)F & OP--F analog & {CCOP(C)(=O)F} \\
GB\_sarin & iPrO--P(Me)(=O)F & OP--F analog & {CC(C)OP(C)(=O)F} \\
OP\_Et\_Et & EtO--P(Et)(=O)F & OP--F analog & {CCOP(=O)(F)CC} \\
OP\_Et\_cHx & EtO--P(cHx)(=O)F & OP--F analog & {CCOP(=O)(F)C1CCCCC1} \\
OP\_Et\_cPn & EtO--P(cPn)(=O)F & OP--F analog & {CCOP(=O)(F)C1CCCC1} \\
OP\_Et\_iBu & EtO--P(iBu)(=O)F & OP--F analog & {CCOP(=O)(F)CC(C)C} \\
OP\_Et\_iPr & EtO--P(iPr)(=O)F & OP--F analog & {CCOP(=O)(F)C(C)C} \\
OP\_Et\_nPr & EtO--P(nPr)(=O)F & OP--F analog & {CCCP(=O)(F)OCC} \\
OP\_Et\_sBu & EtO--P(sBu)(=O)F & OP--F analog & {CCOP(=O)(F)C(C)CC} \\
OP\_Me\_Et & MeO--P(Et)(=O)F & OP--F analog & {CCP(=O)(F)OC} \\
OP\_Me\_Me & MeO--P(Me)(=O)F & OP--F analog & {COP(C)(=O)F} \\
OP\_Me\_cHx & MeO--P(cHx)(=O)F & OP--F analog & {COP(=O)(F)C1CCCCC1} \\
OP\_Me\_cPn & MeO--P(cPn)(=O)F & OP--F analog & {COP(=O)(F)C1CCCC1} \\
OP\_Me\_iBu & MeO--P(iBu)(=O)F & OP--F analog & {COP(=O)(F)CC(C)C} \\
OP\_Me\_iPr & MeO--P(iPr)(=O)F & OP--F analog & {COP(=O)(F)C(C)C} \\
OP\_Me\_nPr & MeO--P(nPr)(=O)F & OP--F analog & {CCCP(=O)(F)OC} \\
OP\_Me\_sBu & MeO--P(sBu)(=O)F & OP--F analog & {CCC(C)P(=O)(F)OC} \\
OP\_cHx\_Et & cHxO--P(Et)(=O)F & OP--F analog & {CCP(=O)(F)OC1CCCCC1} \\
OP\_cPn\_Et & cPnO--P(Et)(=O)F & OP--F analog & {CCP(=O)(F)OC1CCCC1} \\
OP\_cPn\_Me & cPnO--P(Me)(=O)F & OP--F analog & {CP(=O)(F)OC1CCCC1} \\
OP\_cPn\_iPr & cPnO--P(iPr)(=O)F & OP--F analog & {CC(C)P(=O)(F)OC1CCCC1} \\
OP\_cPn\_nPr & cPnO--P(nPr)(=O)F & OP--F analog & {CCCP(=O)(F)OC1CCCC1} \\
OP\_iBu\_Et & iBuO--P(Et)(=O)F & OP--F analog & {CCP(=O)(F)OCC(C)C} \\
OP\_iBu\_iPr & iBuO--P(iPr)(=O)F & OP--F analog & {CC(C)COP(=O)(F)C(C)C} \\
OP\_iBu\_nPr & iBuO--P(nPr)(=O)F & OP--F analog & {CCCP(=O)(F)OCC(C)C} \\
OP\_iPr\_Et & iPrO--P(Et)(=O)F & OP--F analog & {CCP(=O)(F)OC(C)C} \\
OP\_iPr\_cPn & iPrO--P(cPn)(=O)F & OP--F analog & {CC(C)OP(=O)(F)C1CCCC1} \\
OP\_iPr\_iBu & iPrO--P(iBu)(=O)F & OP--F analog & {CC(C)CP(=O)(F)OC(C)C} \\
OP\_iPr\_iPr & iPrO--P(iPr)(=O)F & OP--F analog & {CC(C)OP(=O)(F)C(C)C} \\
OP\_iPr\_nPr & iPrO--P(nPr)(=O)F & OP--F analog & {CCCP(=O)(F)OC(C)C} \\
OP\_iPr\_sBu & iPrO--P(sBu)(=O)F & OP--F analog & {CCC(C)P(=O)(F)OC(C)C} \\
OP\_nPr\_Et & nPrO--P(Et)(=O)F & OP--F analog & {CCCOP(=O)(F)CC} \\
OP\_nPr\_cPn & nPrO--P(cPn)(=O)F & OP--F analog & {CCCOP(=O)(F)C1CCCC1} \\
OP\_nPr\_iBu & nPrO--P(iBu)(=O)F & OP--F analog & {CCCOP(=O)(F)CC(C)C} \\
OP\_nPr\_iPr & nPrO--P(iPr)(=O)F & OP--F analog & {CCCOP(=O)(F)C(C)C} \\
OP\_nPr\_nPr & nPrO--P(nPr)(=O)F & OP--F analog & {CCCOP(=O)(F)CCC} \\
OP\_nPr\_sBu & nPrO--P(sBu)(=O)F & OP--F analog & {CCCOP(=O)(F)C(C)CC} \\
OP\_sBu\_Et & sBuO--P(Et)(=O)F & OP--F analog & {CCC(C)OP(=O)(F)CC} \\
OP\_sBu\_Me & sBuO--P(Me)(=O)F & OP--F analog & {CCC(C)OP(C)(=O)F} \\
OP\_sBu\_iPr & sBuO--P(iPr)(=O)F & OP--F analog & {CCC(C)OP(=O)(F)C(C)C} \\
OP\_sBu\_nPr & sBuO--P(nPr)(=O)F & OP--F analog & {CCCP(=O)(F)OC(C)CC} \\
iBuGB & iBuO--P(Me)(=O)F & OP--F analog & {CC(C)COP(C)(=O)F} \\
nPrGB & nPrO--P(Me)(=O)F & OP--F analog & {CCCOP(C)(=O)F} \\
\end{longtable}}

\begin{flushleft}\footnotesize All spectra computed at B3LYP/6-311++G(d,p) with ORCA 6.1.1.\end{flushleft}
%%%%%%%%%%%%%%%%%%%%%%%%%%%%%%%%%%%%%%%%%%%%%%%%%%%%%%%%%%%%%%%%%%%%%
\clearpage

\subsection{Construction of the analog grid}

The analogs span an $8\times8$ grid of alkyl substituents (methyl, ethyl,
\textit{n}-propyl, \textit{iso}-propyl, \textit{iso}-butyl, \textit{sec}-butyl,
cyclopentyl, cyclohexyl) at R\textsuperscript{1} and R\textsuperscript{2}, restricted to
$\leq$28 atoms, hydrogens included. Table~\ref{tab:si-grid} gives the heavy-atom count of every cell.
With the exception of GF none of them is ever the correct answer; their role is to test
whether identification degrades as the library grows.

%%%%%%%%%%%%%%%%%%%%%%%%%%%%%%%%%%%%%%%%%%%%%%%%%%%%%%%%%%%%%%%%%%%%%
%%  table: grid
%%%%%%%%%%%%%%%%%%%%%%%%%%%%%%%%%%%%%%%%%%%%%%%%%%%%%%%%%%%%%%%%%%%%%
\begin{table}[htbp]
\centering
\caption{Atom count, hydrogens included, of every cell of the $8\times8$ R$^1$O--P(R$^2$)(=O)F substituent grid. Rows are the O-alkyl substituent R$^1$, columns the P-alkyl substituent R$^2$. Cells in bold exceed the 28-atom ceiling and were not computed; the remaining 44 cells constitute the analog set.}
\label{tab:si-grid}
\small
\begin{tabular}{lrrrrrrrr}
\toprule
R$^1$O $\backslash$ R$^2$ & Me & Et & nPr & iPr & iBu & sBu & cPn & cHx \\
\midrule
Me & 12 & 15 & 18 & 18 & 21 & 21 & 22 & 25 \\
Et & 15 & 18 & 21 & 21 & 24 & 24 & 25 & 28 \\
nPr & 18 & 21 & 24 & 24 & 27 & 27 & 28 & \textbf{31} \\
iPr & 18 & 21 & 24 & 24 & 27 & 27 & 28 & \textbf{31} \\
iBu & 21 & 24 & 27 & 27 & \textbf{30} & \textbf{30} & \textbf{31} & \textbf{34} \\
sBu & 21 & 24 & 27 & 27 & \textbf{30} & \textbf{30} & \textbf{31} & \textbf{34} \\
cPn & 22 & 25 & 28 & 28 & \textbf{31} & \textbf{31} & \textbf{32} & \textbf{35} \\
cHx & 25 & 28 & \textbf{31} & \textbf{31} & \textbf{34} & \textbf{34} & \textbf{35} & \textbf{38} \\
\bottomrule
\end{tabular}

\begin{flushleft}\footnotesize Substituents: Me = methyl, Et = ethyl, nPr = $n$-propyl, iPr = $iso$-propyl, iBu = $iso$-butyl, sBu = $sec$-butyl, cPn = cyclopentyl, cHx = cyclohexyl. The ceiling is on the total atom count, which is what sets the cost of a numerical frequency calculation; it was imposed to keep each job tractable. For reference, cyclosarin occupies the cyclohexyl/methyl cell at 25 atoms, of which 11 are heavy.\end{flushleft}
\end{table}
%%%%%%%%%%%%%%%%%%%%%%%%%%%%%%%%%%%%%%%%%%%%%%%%%%%%%%%%%%%%%%%%%%%%%
\clearpage

\subsection{Computational details and vibrational analysis}
\label{si:compdetails}

Starting geometries were generated from SMILES with RDKit (ETKDG embedding followed by
MMFF94 pre-optimization) and written as Cartesian coordinates. Every molecule was then
optimized with ORCA 6.1.1 at the B3LYP/6-311++G(d,p) level using tight convergence
criteria (\texttt{TightOpt}), followed by a numerical frequency calculation
(\texttt{NumFreq}) with polarizability derivatives enabled (\texttt{\%elprop Polar 1}),
which yields harmonic wavenumbers together with Raman activities
$S_i = (\partial\alpha/\partial Q_i)^2$. A representative input header is

\begin{verbatim}
! B3LYP 6-311++G(d,p) TightOpt NumFreq
%pal nprocs 8 end
%maxcore 3000
%elprop Polar 1 end
\end{verbatim}

All 54 structures converged to true minima: none has an imaginary frequency
(Table~\ref{tab:si-freq}). The lowest real wavenumbers, 5--73\,\wn{}, are soft torsional
modes that lie below the 250\,\wn{} lower bound of the analysis window.

Each calculation was run in its own scratch directory and any orphaned MPI processes were
cleared before the next job started; without this, a failure in one job propagated into
the next. We note it because the same failure mode is easy to mistake for a chemistry
problem. The calculations used to test higher levels of theory
(Section~\ref{si:theory}) were run identically except for the keywords stated there.

%%%%%%%%%%%%%%%%%%%%%%%%%%%%%%%%%%%%%%%%%%%%%%%%%%%%%%%%%%%%%%%%%%%%%
%%  table: freq
%%%%%%%%%%%%%%%%%%%%%%%%%%%%%%%%%%%%%%%%%%%%%%%%%%%%%%%%%%%%%%%%%%%%%
{\small
\begin{longtable}{lrrr|lrrr}
\caption{Vibrational analysis of every library member. $N_\mathrm{real}$ is the number of real harmonic modes, $N_\mathrm{imag}$ the number of imaginary modes, and $\nu_\mathrm{min}$ the lowest real wavenumber.}\label{tab:si-freq}\\
\toprule
Job & $N_\mathrm{real}$ & $N_\mathrm{imag}$ & $\nu_\mathrm{min}$ (\wn) & Job & $N_\mathrm{real}$ & $N_\mathrm{imag}$ & $\nu_\mathrm{min}$ (\wn) \\
\midrule
\endfirsthead
\multicolumn{8}{l}{\footnotesize\emph{Table \ref{tab:si-freq}, continued}}\\
\toprule
Job & $N_\mathrm{real}$ & $N_\mathrm{imag}$ & $\nu_\mathrm{min}$ (\wn) & Job & $N_\mathrm{real}$ & $N_\mathrm{imag}$ & $\nu_\mathrm{min}$ (\wn) \\
\midrule
\endhead
\bottomrule
\endfoot
EXP\_A230 & 78 & 0 & 28.1 & OP\_cHx\_Et & 78 & 0 & 22.1 \\
EXP\_A232 & 81 & 0 & 32.2 & OP\_cPn\_Et & 69 & 0 & 25.0 \\
EXP\_A234 & 90 & 0 & 18.1 & OP\_cPn\_Me & 60 & 0 & 30.2 \\
EXP\_A242 & 111 & 0 & 39.2 & OP\_cPn\_iPr & 78 & 0 & 26.9 \\
EXP\_GD\_soman & 75 & 0 & 25.8 & OP\_cPn\_nPr & 78 & 0 & 24.1 \\
EXP\_HD & 39 & 0 & 27.1 & OP\_iBu\_Et & 66 & 0 & 10.8 \\
EXP\_HN1 & 60 & 0 & 27.4 & OP\_iBu\_iPr & 75 & 0 & 27.0 \\
EXP\_HN2 & 51 & 0 & 45.0 & OP\_iBu\_nPr & 75 & 0 & 15.3 \\
EXP\_HN3 & 60 & 0 & 23.4 & OP\_iPr\_Et & 57 & 0 & 23.3 \\
EtGB & 39 & 0 & 43.2 & OP\_iPr\_cPn & 78 & 0 & 18.5 \\
GB\_sarin & 48 & 0 & 22.4 & OP\_iPr\_iBu & 75 & 0 & 19.6 \\
GF\_cyclosarin & 69 & 0 & 25.5 & OP\_iPr\_iPr & 66 & 0 & 15.7 \\
OP\_Et\_Et & 48 & 0 & 32.4 & OP\_iPr\_nPr & 66 & 0 & 33.4 \\
OP\_Et\_cHx & 78 & 0 & 29.3 & OP\_iPr\_sBu & 75 & 0 & 25.5 \\
OP\_Et\_cPn & 69 & 0 & 24.1 & OP\_nPr\_Et & 57 & 0 & 28.0 \\
OP\_Et\_iBu & 66 & 0 & 13.7 & OP\_nPr\_cPn & 78 & 0 & 26.0 \\
OP\_Et\_iPr & 57 & 0 & 29.5 & OP\_nPr\_iBu & 75 & 0 & 10.4 \\
OP\_Et\_nPr & 57 & 0 & 45.0 & OP\_nPr\_iPr & 66 & 0 & 21.2 \\
OP\_Et\_sBu & 66 & 0 & 31.6 & OP\_nPr\_nPr & 66 & 0 & 17.1 \\
OP\_Me\_Et & 39 & 0 & 72.6 & OP\_nPr\_sBu & 75 & 0 & 26.5 \\
OP\_Me\_Me & 30 & 0 & 72.2 & OP\_sBu\_Et & 66 & 0 & 24.2 \\
OP\_Me\_cHx & 69 & 0 & 30.7 & OP\_sBu\_Me & 57 & 0 & 8.8 \\
OP\_Me\_cPn & 60 & 0 & 30.3 & OP\_sBu\_iPr & 75 & 0 & 25.4 \\
OP\_Me\_iBu & 57 & 0 & 35.1 & OP\_sBu\_nPr & 75 & 0 & 24.1 \\
OP\_Me\_iPr & 48 & 0 & 48.0 & iBuGB & 57 & 0 & 4.9 \\
OP\_Me\_nPr & 48 & 0 & 57.5 & nPrGB & 48 & 0 & 30.6 \\
OP\_Me\_sBu & 57 & 0 & 31.3 &  &  &  &  \\
\end{longtable}}

\begin{flushleft}\footnotesize No structure in the library has an imaginary frequency, so all 54 geometries are true minima on the B3LYP/6-311++G(d,p) surface. The lowest real wavenumbers (5--73\,\wn) are soft torsions that fall below the 250\,\wn{} lower bound of the analysis window and therefore do not enter the spectral comparison.\end{flushleft}
%%%%%%%%%%%%%%%%%%%%%%%%%%%%%%%%%%%%%%%%%%%%%%%%%%%%%%%%%%%%%%%%%%%%%
\clearpage

%%%%%%%%%%%%%%%%%%%%%%%%%%%%%%%%%%%%%%%%%%%%%%%%%%%%%%%%%%%%%%%%%%%%%
\section{Calibration}
\label{si:calib}

\subsection{The linewidth is measured}

For each acquisition a single Lorentzian width shared by all detected bands was fitted
simultaneously with the band amplitudes, band centers being free to move by at most one
pixel. Table~\ref{tab:si-gamma} gives the result per agent. We adopt
$\Gamma = 22$\,\wn{}, which was obtained from the measured spectra alone.

%%%%%%%%%%%%%%%%%%%%%%%%%%%%%%%%%%%%%%%%%%%%%%%%%%%%%%%%%%%%%%%%%%%%%
%%  table: gamma
%%%%%%%%%%%%%%%%%%%%%%%%%%%%%%%%%%%%%%%%%%%%%%%%%%%%%%%%%%%%%%%%%%%%%
\begin{table}[htbp]
\centering
\caption{The Lorentzian linewidth is measured, not fitted to identification. For each acquisition a single width shared by all detected bands was fitted simultaneously with the band amplitudes, band centers being free to move by at most one pixel.}
\label{tab:si-gamma}
\small
\begin{tabular}{lrrrrr}
\toprule
Agent & Acquisitions & $\Gamma$ median (\wn) & Range (\wn) & $R^2$ median & Bands fitted \\
\midrule
GD & 4 & 18.9 & 18.4--19.8 & 0.92 & 19 \\
GF & 3 & 17.4 & 17.4--17.7 & 0.91 & 20 \\
VX & 1 & 38.9 & 38.9--38.9 & 0.89 & 17 \\
HD & 1 & 48.4 & 48.4--48.4 & 0.81 & 7 \\
HN1 & 3 & 21.7 & 21.6--21.7 & 0.91 & 16 \\
HN2 & 3 & 9.3 & 9.1--9.3 & 0.98 & 7 \\
HN3 & 3 & 22.5 & 22.3--22.7 & 0.87 & 6 \\
A-230 & 1 & 17.8 & 17.8--17.8 & 0.82 & 10 \\
A-232 & 3 & 24.4 & 24.2--25.5 & 0.92 & 18 \\
A-234 & 4 & 29.8 & 21.6--30.0 & 0.91 & 17 \\
A-242 & 2 & 33.7 & 32.2--35.2 & 0.91 & 18 \\
\bottomrule
\end{tabular}

\begin{flushleft}\footnotesize Of the 32 genuine acquisitions, 28 yielded a converged fit ($R^2 > 0.8$); the remaining 4 did not and are excluded from this table. Over the 28 the median is 21.7\,\wn{} with an interquartile range of 18.2--26.5\,\wn{}. We adopt $\Gamma$ = 22\,\wn{}. The extremes are not trustworthy: HN2 offers a single band, and HD and VX a single acquisition each. This is an apparent width that includes the instrument response function. No identification result was consulted in obtaining it.\end{flushleft}
\end{table}
%%%%%%%%%%%%%%%%%%%%%%%%%%%%%%%%%%%%%%%%%%%%%%%%%%%%%%%%%%%%%%%%%%%%%
\clearpage

\subsection{Band-position agreement}
\label{si:bands}

The frequency scale was fitted by minimizing the root-mean-square difference between
measured band maxima and scaled harmonic wavenumbers, and Table~\ref{tab:si-bands}
reports the agreement it achieves. Table~\ref{tab:si-bandlist} lists every matched pair.
No coefficient of determination is reported, for the reason given in
Section~\ref{si:corrections}.

%%%%%%%%%%%%%%%%%%%%%%%%%%%%%%%%%%%%%%%%%%%%%%%%%%%%%%%%%%%%%%%%%%%%%
%%  table: bands
%%%%%%%%%%%%%%%%%%%%%%%%%%%%%%%%%%%%%%%%%%%%%%%%%%%%%%%%%%%%%%%%%%%%%
\begin{table}[htbp]
\centering
\caption{Agreement between calibrated computed and measured band positions, the quantity the frequency scale was fitted to. Measured and computed maxima are paired by one-to-one assignment, solved as a linear assignment problem on the matrix of absolute wavenumber differences with pairings beyond 35\,\wn{} forbidden; computed bands below 2\% of the molecule's strongest activity are not considered. Nearest-neighbor matching is not used, because it allows one computed band to be claimed by several measured bands with residuals of opposite sign that cancel.}
\label{tab:si-bands}
\small
\begin{tabular}{lrrrrr}
\toprule
Agent & Pairs & RMSE (\wn) & Mean residual & Min & Max \\
\midrule
GD & 8 & 7.5 & -0.2 & -8 & +16 \\
GF & 4 & 11.4 & -8.6 & -21 & -1 \\
VX & 8 & 9.1 & -2.8 & -15 & +11 \\
HD & 11 & 8.0 & -2.7 & -15 & +10 \\
HN1 & 6 & 9.5 & +4.3 & -9 & +14 \\
HN3 & 7 & 11.8 & -2.5 & -23 & +9 \\
A-230 & 13 & 10.8 & +2.5 & -11 & +29 \\
A-232 & 8 & 11.4 & +0.7 & -19 & +19 \\
A-234 & 7 & 15.6 & +7.3 & -22 & +28 \\
A-242 & 9 & 8.5 & -3.8 & -12 & +16 \\
\textbf{All} & \textbf{81} & \textbf{10.4} & \textbf{-0.3} & \textbf{-23} & \textbf{+29} \\
\bottomrule
\end{tabular}

\begin{flushleft}\footnotesize The global RMSE of 10.4\,\wn{} is approximately 1.8 pixels on the native measurement grid, whose mean spacing in this region is 5.73\,\wn{}. The full list of pairs is given in Table~\ref{tab:si-bandlist}.\end{flushleft}
\end{table}
%%%%%%%%%%%%%%%%%%%%%%%%%%%%%%%%%%%%%%%%%%%%%%%%%%%%%%%%%%%%%%%%%%%%%
\clearpage

%%%%%%%%%%%%%%%%%%%%%%%%%%%%%%%%%%%%%%%%%%%%%%%%%%%%%%%%%%%%%%%%%%%%%
%%  table: bandlist
%%%%%%%%%%%%%%%%%%%%%%%%%%%%%%%%%%%%%%%%%%%%%%%%%%%%%%%%%%%%%%%%%%%%%
\begin{table}[p]
\centering
\caption{All 81 matched band pairs used for the frequency calibration. Positions in \wn{}; the residual is calibrated computed minus measured.}
\label{tab:si-bandlist}
\footnotesize
\resizebox{0.84\textwidth}{!}{%
\begin{tabular}{lrrr@{\hskip 2em}lrrr}
\toprule
Agent & Meas. & Comp. & $\Delta$ & Agent & Meas. & Comp. & $\Delta$ \\
\midrule
GD & 540 & 543 & +3 & HN3 & 1143 & 1127 & -16 \\
GD & 752 & 745 & -6 & HN3 & 1232 & 1209 & -23 \\
GD & 869 & 861 & -7 & HN3 & 1444 & 1453 & +9 \\
GD & 935 & 927 & -8 & A-230 & 553 & 554 & +0 \\
GD & 1226 & 1228 & +2 & A-230 & 708 & 721 & +14 \\
GD & 1264 & 1266 & +2 & A-230 & 739 & 728 & -11 \\
GD & 1355 & 1372 & +16 & A-230 & 976 & 966 & -10 \\
GD & 1449 & 1447 & -3 & A-230 & 1046 & 1075 & +29 \\
GF & 832 & 811 & -21 & A-230 & 1080 & 1089 & +9 \\
GF & 1023 & 1017 & -6 & A-230 & 1215 & 1218 & +3 \\
GF & 1270 & 1262 & -7 & A-230 & 1243 & 1241 & -1 \\
GF & 1444 & 1444 & -1 & A-230 & 1291 & 1297 & +6 \\
VX & 460 & 454 & -6 & A-230 & 1377 & 1367 & -10 \\
VX & 695 & 691 & -4 & A-230 & 1454 & 1451 & -3 \\
VX & 739 & 741 & +1 & A-230 & 1500 & 1508 & +7 \\
VX & 887 & 878 & -9 & A-230 & 1600 & 1601 & +1 \\
VX & 1098 & 1109 & +11 & A-232 & 580 & 561 & -19 \\
VX & 1215 & 1205 & -10 & A-232 & 708 & 727 & +19 \\
VX & 1307 & 1293 & -15 & A-232 & 758 & 755 & -3 \\
VX & 1460 & 1469 & +9 & A-232 & 976 & 966 & -10 \\
HD & 644 & 645 & +1 & A-232 & 1080 & 1088 & +8 \\
HD & 702 & 696 & -6 & A-232 & 1291 & 1302 & +11 \\
HD & 733 & 740 & +7 & A-232 & 1454 & 1451 & -4 \\
HD & 758 & 768 & +10 & A-232 & 1495 & 1499 & +3 \\
HD & 869 & 858 & -10 & A-234 & 714 & 723 & +9 \\
HD & 1034 & 1033 & -1 & A-234 & 758 & 766 & +8 \\
HD & 1204 & 1210 & +5 & A-234 & 999 & 977 & -22 \\
HD & 1291 & 1288 & -3 & A-234 & 1080 & 1108 & +28 \\
HD & 1324 & 1313 & -10 & A-234 & 1291 & 1304 & +12 \\
HD & 1439 & 1432 & -7 & A-234 & 1454 & 1462 & +8 \\
HD & 1485 & 1470 & -15 & A-234 & 1500 & 1509 & +8 \\
HN1 & 663 & 655 & -9 & A-242 & 702 & 693 & -8 \\
HN1 & 720 & 716 & -5 & A-242 & 733 & 726 & -7 \\
HN1 & 770 & 778 & +7 & A-242 & 941 & 936 & -4 \\
HN1 & 1040 & 1053 & +13 & A-242 & 999 & 992 & -7 \\
HN1 & 1297 & 1311 & +14 & A-242 & 1075 & 1091 & +16 \\
HN1 & 1449 & 1455 & +5 & A-242 & 1215 & 1215 & -0 \\
HN3 & 663 & 658 & -5 & A-242 & 1253 & 1250 & -4 \\
HN3 & 720 & 728 & +8 & A-242 & 1307 & 1295 & -12 \\
HN3 & 946 & 954 & +7 & A-242 & 1454 & 1447 & -7 \\
HN3 & 1011 & 1013 & +2 &  &  &  &  \\
\bottomrule
\end{tabular}
}
\end{table}
%%%%%%%%%%%%%%%%%%%%%%%%%%%%%%%%%%%%%%%%%%%%%%%%%%%%%%%%%%%%%%%%%%%%%
\clearpage

Two qualifications belong with that number. The pairing is one-to-one, solved as a linear
assignment problem, because nearest-neighbor matching allows one computed band to be
claimed by several measured bands with residuals of opposite sign that cancel; the previous
version of this analysis used nearest-neighbor matching and is corrected here
(Section~\ref{si:corrections}). And the root-mean-square difference is taken over the bands
that found a partner within the tolerance, which by construction excludes the worst-matched
bands. Table~\ref{tab:si-census} therefore reports the complete census, including the bands
that found no partner in either direction, and charges each of them the full tolerance as a
penalty so that the optimism of the gated figure is visible.

%%%%%%%%%%%%%%%%%%%%%%%%%%%%%%%%%%%%%%%%%%%%%%%%%%%%%%%%%%%%%%%%%%%%%
%%  table: census
%%%%%%%%%%%%%%%%%%%%%%%%%%%%%%%%%%%%%%%%%%%%%%%%%%%%%%%%%%%%%%%%%%%%%
\begin{table}[htbp]
\centering
\caption{Complete pairing census for the frequency calibration. Bands that find no partner within the 35\,\wn{} tolerance are counted rather than discarded silently.}
\label{tab:si-census}
\small
\begin{tabular}{lrrrrr}
\toprule
Agent & Measured bands & Computed bands & Paired & Unpaired meas. & Unpaired comp. \\
\midrule
GD & 13 & 15 & 8 & 5 & 7 \\
GF & 11 & 13 & 4 & 7 & 9 \\
VX & 15 & 13 & 8 & 7 & 5 \\
HD & 40 & 14 & 11 & 29 & 3 \\
HN1 & 7 & 19 & 6 & 1 & 13 \\
HN2 & 3 & 10 & 0 & 3 & 10 \\
HN3 & 10 & 24 & 7 & 3 & 17 \\
A-230 & 29 & 22 & 13 & 16 & 9 \\
A-232 & 14 & 19 & 8 & 6 & 11 \\
A-234 & 12 & 8 & 7 & 5 & 1 \\
A-242 & 22 & 18 & 9 & 13 & 9 \\
\textbf{All} & \textbf{176} & \textbf{175} & \textbf{81} & \textbf{95} & \textbf{94} \\
\bottomrule
\end{tabular}

\begin{flushleft}\footnotesize The positional RMSE quoted in the text is computed over the 81 paired bands only, which by construction excludes the worst-matched bands and is therefore optimistic. Charging every unpaired band the full tolerance as a penalty gives 29.8\,\wn{} instead of 10.4\,\wn{}. Both are reported so that the effect of the tolerance gate is visible.\end{flushleft}
\end{table}
%%%%%%%%%%%%%%%%%%%%%%%%%%%%%%%%%%%%%%%%%%%%%%%%%%%%%%%%%%%%%%%%%%%%%
\clearpage

Table~\ref{tab:si-detect} varies the band-detection criterion and the pairing tolerance
together. Neither drives the calibration: the two scale factors move by at most 0.2\% and
the paired RMSE stays between 8.2 and 10.6\,\wn{} across three criteria and two
tolerances.

%%%%%%%%%%%%%%%%%%%%%%%%%%%%%%%%%%%%%%%%%%%%%%%%%%%%%%%%%%%%%%%%%%%%%
%%  table: detect
%%%%%%%%%%%%%%%%%%%%%%%%%%%%%%%%%%%%%%%%%%%%%%%%%%%%%%%%%%%%%%%%%%%%%
\begin{table}[htbp]
\centering
\caption{Sensitivity of the frequency calibration to the band-detection criterion and to the pairing tolerance. The criterion in the first row is the one used throughout this work.}
\label{tab:si-detect}
\footnotesize
\resizebox{\textwidth}{!}{%
\begin{tabular}{lrrrrrrrr}
\toprule
Detection criterion & Tol.\ (\wn) & $s_\mathrm{low}$ & $s_\mathrm{high}$ & RMSE (\wn) & Paired & Meas.\ bands & Comp.\ bands & Penalized RMSE \\
\midrule
prominence $>$ 0.05, separation 2 points & 35 & 1.016 & 0.990 & 10.4 & 81 & 176 & 175 & 29.8 \\
prominence $>$ 0.05, separation 2 points & 20 & 1.016 & 0.990 & 8.6 & 76 & 176 & 175 & 30.1 \\
height $>$ 0.10 and prominence $>$ 0.05, separation 4 points & 35 & 1.016 & 0.988 & 10.6 & 75 & 132 & 175 & 29.4 \\
height $>$ 0.10 and prominence $>$ 0.05, separation 4 points & 20 & 1.016 & 0.988 & 8.3 & 69 & 132 & 175 & 29.8 \\
height $>$ 0.12, separation 6 points & 35 & 1.016 & 0.988 & 10.4 & 71 & 121 & 175 & 29.5 \\
height $>$ 0.12, separation 6 points & 20 & 1.016 & 0.988 & 8.2 & 65 & 121 & 175 & 30.0 \\
\bottomrule
\end{tabular}
}

\begin{flushleft}\footnotesize The scale factors move by at most 0.2\% and the paired RMSE stays between 8.2 and 10.6\,\wn{} across all six combinations, so neither the detection threshold nor the tolerance drives the calibration. The penalized RMSE charges every unpaired band a fixed 35\,\wn{} so that the six cells remain comparable; it is stable near 30\,\wn{} and is the honest upper bound on the positional disagreement, the paired RMSE being the value over bands that found a partner.\end{flushleft}
\end{table}
%%%%%%%%%%%%%%%%%%%%%%%%%%%%%%%%%%%%%%%%%%%%%%%%%%%%%%%%%%%%%%%%%%%%%
\clearpage


\subsection{Held-out calibration}

The calibration is fitted on band positions from all agents, so a reader may reasonably
ask whether an agent's own bands carry it. Table~\ref{tab:si-loao} refits the two scale
factors with each agent held out in turn.

%%%%%%%%%%%%%%%%%%%%%%%%%%%%%%%%%%%%%%%%%%%%%%%%%%%%%%%%%%%%%%%%%%%%%
%%  table: loao
%%%%%%%%%%%%%%%%%%%%%%%%%%%%%%%%%%%%%%%%%%%%%%%%%%%%%%%%%%%%%%%%%%%%%
\begin{table}[htbp]
\centering
\caption{Nested leave-one-agent-out calibration. The two frequency scale factors were refitted with each agent's band positions excluded in turn.}
\label{tab:si-loao}
\small
\begin{tabular}{lrrrr}
\toprule
Agent held out & $s_\mathrm{low}$ & $s_\mathrm{high}$ & RMSE (\wn) & Band pairs \\
\midrule
GD & 1.0160 & 0.9960 & 10.64 & 72 \\
GF & 1.0160 & 0.9900 & 10.35 & 77 \\
VX & 1.0160 & 0.9900 & 10.54 & 73 \\
HD & 1.0160 & 0.9880 & 10.68 & 70 \\
HN1 & 1.0160 & 0.9960 & 10.21 & 74 \\
HN2 & 1.0160 & 0.9900 & 10.41 & 81 \\
HN3 & 1.0160 & 0.9900 & 10.26 & 74 \\
A-230 & 1.0160 & 0.9900 & 10.33 & 68 \\
A-232 & 1.0160 & 0.9900 & 10.29 & 73 \\
A-234 & 1.0160 & 0.9960 & 9.52 & 73 \\
A-242 & 1.0140 & 0.9900 & 10.61 & 72 \\
\textbf{All agents} & \textbf{1.0160} & \textbf{0.9900} & \textbf{10.41} & \textbf{81} \\
\bottomrule
\end{tabular}

\begin{flushleft}\footnotesize $s_\mathrm{low}$ varies over 1.0140--1.0160, a range of 0.20\%, and $s_\mathrm{high}$ over 0.9880--0.9960, a range of 0.81\%. No agent drives the calibration, and the identification results are obtained under factors effectively independent of the agent being identified.\end{flushleft}
\end{table}
%%%%%%%%%%%%%%%%%%%%%%%%%%%%%%%%%%%%%%%%%%%%%%%%%%%%%%%%%%%%%%%%%%%%%
\clearpage

%%%%%%%%%%%%%%%%%%%%%%%%%%%%%%%%%%%%%%%%%%%%%%%%%%%%%%%%%%%%%%%%%%%%%
\begin{table}[htbp]
\centering
\caption{Identification recomputed under each held-out calibration. For each agent the
two scale factors are refitted with that agent's band pairs removed, and the whole
identification is then rerun with those factors. This is the table the claim in the main
text rests on.}
\label{tab:si-loaoident}
\small
\begin{tabular}{lrrrcc}
\toprule
Agent held out & $s_\mathrm{low}$ & $s_\mathrm{high}$ & RMSE (\wn) & Molecular top-1 & Class \\
\midrule
--- (all agents) & 1.016 & 0.990 & 10.41 & \textbf{13/30} & 25/30 \\
\midrule
GD & 1.016 & 0.996 & 10.64 & 13/30 & 22/30 \\
GF & 1.016 & 0.990 & 10.35 & 13/30 & 25/30 \\
VX & 1.016 & 0.990 & 10.54 & 13/30 & 25/30 \\
HD & 1.016 & 0.988 & 10.68 & 13/30 & 25/30 \\
HN1 & 1.016 & 0.996 & 10.21 & 13/30 & 22/30 \\
HN2 & 1.016 & 0.990 & 10.41 & 13/30 & 25/30 \\
HN3 & 1.016 & 0.990 & 10.26 & 13/30 & 25/30 \\
A-230 & 1.016 & 0.990 & 10.33 & 13/30 & 25/30 \\
A-232 & 1.016 & 0.990 & 10.29 & 13/30 & 25/30 \\
A-234 & 1.016 & 0.996 & 9.52 & 13/30 & 22/30 \\
A-242 & 1.014 & 0.990 & 10.61 & 13/30 & 25/30 \\
\bottomrule
\end{tabular}

\begin{flushleft}\footnotesize The molecular result is 13 of 30 under every one of the
eleven held-out calibrations, so the figure the paper reports does not depend on any one
agent's band positions entering the fit. The class-level result is not equally stable:
the three calibrations that push $s_\mathrm{high}$ to 0.996 give 22 of 30 rather than 25,
because two A-234 acquisitions and one further acquisition cross a class boundary. We
report the molecular figure as robust to the calibration and the class figure as
not.\end{flushleft}
\end{table}

\subsection{The wavenumber axis: an affine alternative and an offset scan}
\label{si:axis}

A multiplicative frequency scale factor is degenerate with a multiplicative error in the
instrument's wavenumber axis, and an \emph{additive} axis offset is not excluded by the
two-region fit. An additive offset is the error that would most favour this study's
conclusion, because it cancels in a measured-to-measured comparison and not in a
measured-to-computed one. This section reports the two tests we ran on it.

\begin{table}[htbp]
\centering
\caption{Three calibration models fitted to the same 81 band pairs with the same
non-circular objective (band positions only). The affine model has the same number of
free parameters as the adopted two-region scaling and fits better.}
\label{tab:si-affine}
\small
\begin{tabular}{llr}
\toprule
Model & Fitted parameters & Position RMSE (\wn) \\
\midrule
Single multiplicative scale & $s$ = 0.993 & 14.3 \\
Two-region scale (adopted) & $s_\mathrm{low}$ = 1.016, $s_\mathrm{high}$ = 0.990 at 800\,\wn{} & 10.4 \\
Affine, $g\nu + \delta$ & $g$ = 0.966, $\delta$ = $+$34\,\wn{} & \textbf{8.9} \\
\bottomrule
\end{tabular}

\begin{flushleft}\footnotesize The affine gain of 0.966 is the conventional B3LYP
fundamental scale factor, whereas the adopted $s_\mathrm{low}$ exceeds unity. The two
two-parameter models are not distinguishable as explanations: their predicted band
positions differ by 5.1\,\wn{} RMS and at most 10.8\,\wn{}, against a residual of 8.9--10.4\,\wn{}
for either. Below 800\,\wn{} they cannot be separated at all --- 25 of the 81 pairs lie
there, the design matrix has condition number 5.8 $\times$ 10$^{3}$, and the offset and
slope estimates correlate at $-$0.993.\end{flushleft}
\end{table}

\begin{table}[htbp]
\centering
\caption{Identification under an imposed additive error on the measured wavenumber axis.
At each $\delta$ the measured axis is shifted, the calibration is refitted by the same
non-circular procedure, and the whole identification is recomputed.}
\label{tab:si-offsetscan}
\small
\begin{tabular}{rrrrrrrr}
\toprule
$\delta$ (\wn) & $s_\mathrm{low}$ & $s_\mathrm{high}$ & RMSE (\wn) & Pairs & Computed & Measured & Computed self \\
\midrule
$-$30 & 0.980 & 1.014 & 12.55 & 71 & 11/30 & 30/30 & 0.577 \\
$-$20 & 1.040 & 1.006 & 11.15 & 81 & 13/30 & 30/30 & 0.659 \\
$-$10 & 1.030 & 0.998 & 10.66 & 81 & 13/30 & 30/30 & 0.668 \\
$-$5 & 1.022 & 0.994 & 10.55 & 81 & 13/30 & 30/30 & 0.671 \\
\textbf{0 (adopted)} & \textbf{1.016} & \textbf{0.990} & \textbf{10.41} & \textbf{81} & \textbf{13/30} & \textbf{30/30} & \textbf{0.674} \\
$+$5 & 1.008 & 0.992 & 9.71 & 79 & 13/30 & 30/30 & 0.649 \\
$+$10 & 1.002 & 0.988 & 9.60 & 79 & 13/30 & 30/30 & 0.650 \\
$+$20 & 0.986 & 0.980 & 9.45 & 79 & 13/30 & 30/30 & 0.649 \\
$+$30 & 0.972 & 0.970 & 9.02 & 79 & 13/30 & 30/30 & 0.654 \\
\bottomrule
\end{tabular}

\begin{flushleft}\footnotesize The computed-library result is flat at 13 of 30 over an
imposed axis error of $-$20 to $+$30\,\wn{}, five times the pixel spacing, and the
similarity to the correct entry varies only between 0.649 and 0.674. The fitted scale
factors absorb the offset; the identification does not move. Rebuilding the library under
the affine calibration of Table~\ref{tab:si-affine} likewise gives 13 of 30 with a
self-similarity of 0.666. A wavenumber-standard measurement remains worth having for any
replication, but an uncalibrated axis does not account for the result reported here.\end{flushleft}
\end{table}

%%%%%%%%%%%%%%%%%%%%%%%%%%%%%%%%%%%%%%%%%%%%%%%%%%%%%%%%%%%%%%%%%%%%%
\clearpage

\section{Identification}
\label{si:ident}

\subsection{The same queries against both libraries}

Table~\ref{tab:si-ident} summarizes by agent and Table~\ref{tab:si-peracq} gives every
acquisition. The query acquisition is never present in the library it is matched against.

%%%%%%%%%%%%%%%%%%%%%%%%%%%%%%%%%%%%%%%%%%%%%%%%%%%%%%%%%%%%%%%%%%%%%
%%  table: ident
%%%%%%%%%%%%%%%%%%%%%%%%%%%%%%%%%%%%%%%%%%%%%%%%%%%%%%%%%%%%%%%%%%%%%
\begin{table}[htbp]
\centering
\caption{Identification of the same thirty genuine acquisitions against a measured and against a computed library, under identical preprocessing, an identical native grid and an identical cosine metric. The query acquisition is never present in the library it is matched against.}
\label{tab:si-ident}
\footnotesize
\resizebox{\textwidth}{!}{%
\begin{tabular}{lrrrrrrr}
\toprule
Agent & $n$ & Measured mol. & Self cos. & Computed mol. & Self cos. & Computed class & 54-member mol. \\
\midrule
GD & 4 & 4/4 & 0.999 & 4/4 & 0.853 & 4/4 & 4/4 \\
GF & 3 & 3/3 & 1.000 & 0/3 & 0.517 & 3/3 & 0/3 \\
HN1 & 3 & 3/3 & 1.000 & 0/3 & 0.831 & 3/3 & 0/3 \\
HN2 & 3 & 3/3 & 1.000 & 0/3 & 0.057 & 0/3 & 0/3 \\
HN3 & 3 & 3/3 & 0.999 & 3/3 & 0.830 & 3/3 & 3/3 \\
A-230 & 3 & 3/3 & 0.970 & 3/3 & 0.801 & 3/3 & 3/3 \\
A-232 & 3 & 3/3 & 1.000 & 0/3 & 0.652 & 3/3 & 0/3 \\
A-234 & 5 & 5/5 & 0.941 & 0/5 & 0.722 & 3/5 & 0/5 \\
A-242 & 3 & 3/3 & 0.958 & 3/3 & 0.706 & 3/3 & 3/3 \\
\textbf{All} & \textbf{30} & \textbf{30/30} & \textbf{0.983} & \textbf{13/30} & \textbf{0.674} & \textbf{25/30} & \textbf{13/30} \\
\bottomrule
\end{tabular}
}

\begin{flushleft}\footnotesize HD and VX have a single genuine acquisition each and are retained in both libraries but not queried. The three HN2 queries fail against the computed library at a self-similarity of 0.057; Section~S6 treats that as evidence about the sample.\end{flushleft}
\end{table}
%%%%%%%%%%%%%%%%%%%%%%%%%%%%%%%%%%%%%%%%%%%%%%%%%%%%%%%%%%%%%%%%%%%%%
\clearpage
%%%%%%%%%%%%%%%%%%%%%%%%%%%%%%%%%%%%%%%%%%%%%%%%%%%%%%%%%%%%%%%%%%%%%
%%  table: peracq
%%%%%%%%%%%%%%%%%%%%%%%%%%%%%%%%%%%%%%%%%%%%%%%%%%%%%%%%%%%%%%%%%%%%%
\begin{table}[htbp]
\centering
\caption{Assignment of every individual acquisition. Each row is one genuine acquisition queried against each library in turn. In the OK columns \checkmark{} marks a correct molecular assignment and $\times$ an incorrect one.}
\label{tab:si-peracq}
\footnotesize
\resizebox{\textwidth}{!}{%
\begin{tabular}{llllrrllrr}
\toprule
Agent & Acq. & Measured top-1 & OK & Self & Margin & Computed top-1 & OK & Self & Margin \\
\midrule
A-230 & 1 & A230 & \checkmark & 0.968 & 0.130 & EXP\_A230 & \checkmark & 0.780 & 0.035 \\
A-230 & 2 & A230 & \checkmark & 0.971 & 0.124 & EXP\_A230 & \checkmark & 0.816 & 0.056 \\
A-230 & 3 & A230 & \checkmark & 0.972 & 0.148 & EXP\_A230 & \checkmark & 0.808 & 0.052 \\
A-232 & 1 & A232 & \checkmark & 1.000 & 0.120 & EXP\_A242 & $\times$ & 0.648 & 0.036 \\
A-232 & 2 & A232 & \checkmark & 1.000 & 0.119 & EXP\_A242 & $\times$ & 0.652 & 0.033 \\
A-232 & 3 & A232 & \checkmark & 1.000 & 0.117 & EXP\_A242 & $\times$ & 0.655 & 0.035 \\
A-234 & 1 & A234 & \checkmark & 0.842 & 0.022 & GF\_cyclosarin\_glob & $\times$ & 0.680 & 0.008 \\
A-234 & 2 & A234 & \checkmark & 0.952 & 0.095 & GF\_cyclosarin\_glob & $\times$ & 0.717 & 0.024 \\
A-234 & 3 & A234 & \checkmark & 0.971 & 0.102 & EXP\_A232\_Eglob & $\times$ & 0.737 & 0.000 \\
A-234 & 4 & A234 & \checkmark & 0.971 & 0.103 & EXP\_A232\_Eglob & $\times$ & 0.741 & 0.000 \\
A-234 & 5 & A234 & \checkmark & 0.971 & 0.103 & EXP\_A232\_Eglob & $\times$ & 0.737 & 0.000 \\
A-242 & 1 & A242 & \checkmark & 0.918 & 0.206 & EXP\_A242 & \checkmark & 0.744 & 0.091 \\
A-242 & 2 & A242 & \checkmark & 0.982 & 0.288 & EXP\_A242 & \checkmark & 0.684 & 0.125 \\
A-242 & 3 & A242 & \checkmark & 0.972 & 0.289 & EXP\_A242 & \checkmark & 0.691 & 0.137 \\
GD & 1 & GD & \checkmark & 0.998 & 0.210 & EXP\_GD\_soman & \checkmark & 0.852 & 0.099 \\
GD & 2 & GD & \checkmark & 1.000 & 0.210 & EXP\_GD\_soman & \checkmark & 0.853 & 0.098 \\
GD & 3 & GD & \checkmark & 1.000 & 0.210 & EXP\_GD\_soman & \checkmark & 0.852 & 0.097 \\
GD & 4 & GD & \checkmark & 1.000 & 0.210 & EXP\_GD\_soman & \checkmark & 0.853 & 0.098 \\
GF & 1 & GF & \checkmark & 1.000 & 0.291 & EXP\_GD\_soman & $\times$ & 0.515 & 0.083 \\
GF & 2 & GF & \checkmark & 1.000 & 0.289 & EXP\_GD\_soman & $\times$ & 0.517 & 0.082 \\
GF & 3 & GF & \checkmark & 1.000 & 0.288 & EXP\_GD\_soman & $\times$ & 0.518 & 0.083 \\
HN1 & 1 & HN1 & \checkmark & 1.000 & 0.065 & EXP\_HN3 & $\times$ & 0.832 & 0.004 \\
HN1 & 2 & HN1 & \checkmark & 1.000 & 0.065 & EXP\_HN3 & $\times$ & 0.831 & 0.004 \\
HN1 & 3 & HN1 & \checkmark & 1.000 & 0.065 & EXP\_HN3 & $\times$ & 0.831 & 0.004 \\
HN2 & 1 & HN2 & \checkmark & 1.000 & 0.832 & EXP\_A242 & $\times$ & 0.057 & 0.090 \\
HN2 & 2 & HN2 & \checkmark & 1.000 & 0.833 & EXP\_A242 & $\times$ & 0.058 & 0.090 \\
HN2 & 3 & HN2 & \checkmark & 1.000 & 0.830 & EXP\_A242 & $\times$ & 0.057 & 0.092 \\
HN3 & 1 & HN3 & \checkmark & 0.999 & 0.062 & EXP\_HN3 & \checkmark & 0.833 & 0.037 \\
HN3 & 2 & HN3 & \checkmark & 0.999 & 0.066 & EXP\_HN3 & \checkmark & 0.830 & 0.038 \\
HN3 & 3 & HN3 & \checkmark & 0.999 & 0.065 & EXP\_HN3 & \checkmark & 0.828 & 0.034 \\
\bottomrule
\end{tabular}
}

\begin{flushleft}\footnotesize The agent-mean figures quoted elsewhere in the literature average over these rows. We report the per-acquisition result because it is what a single interrogation delivers.\end{flushleft}
\end{table}
%%%%%%%%%%%%%%%%%%%%%%%%%%%%%%%%%%%%%%%%%%%%%%%%%%%%%%%%%%%%%%%%%%%%%
\clearpage

The thirty acquisitions are repeated interrogations of nine sealed vials recorded within a
single hour, and are not independent sample preparations. Table~\ref{tab:si-agentstat}
therefore repeats the comparison with the agent as the unit of analysis and attaches
interval estimates, and Table~\ref{tab:si-hn2} reports every headline quantity with and
without the HN2 vial, whose contents are not established.

%%%%%%%%%%%%%%%%%%%%%%%%%%%%%%%%%%%%%%%%%%%%%%%%%%%%%%%%%%%%%%%%%%%%%
%%  table: agentstat
%%%%%%%%%%%%%%%%%%%%%%%%%%%%%%%%%%%%%%%%%%%%%%%%%%%%%%%%%%%%%%%%%%%%%
\begin{table}[htbp]
\centering
\caption{Accuracy with the agent, not the acquisition, as the unit of analysis. Each agent is scored by the majority of its acquisitions; intervals are Clopper--Pearson.}
\label{tab:si-agentstat}
\footnotesize
\resizebox{\textwidth}{!}{%
\begin{tabular}{lllc}
\toprule
Subset & Quantity & Agents correct & 95\% CI \\
\midrule
All nine agents & Measured library, molecular & 9/9 & [0.66, 1.00] \\
All nine agents & Computed library, molecular (11-member) & 4/9 & [0.14, 0.79] \\
All nine agents & Computed library, molecular (54-member) & 4/9 & [0.14, 0.79] \\
All nine agents & Measured library, class & 9/9 & [0.66, 1.00] \\
All nine agents & Computed library, class & 8/9 & [0.52, 1.00] \\
HN2 excluded (eight agents) & Measured library, molecular & 8/8 & [0.63, 1.00] \\
HN2 excluded (eight agents) & Computed library, molecular (11-member) & 4/8 & [0.16, 0.84] \\
HN2 excluded (eight agents) & Computed library, molecular (54-member) & 4/8 & [0.16, 0.84] \\
HN2 excluded (eight agents) & Measured library, class & 8/8 & [0.63, 1.00] \\
HN2 excluded (eight agents) & Computed library, class & 8/8 & [0.63, 1.00] \\
\bottomrule
\end{tabular}
}

\begin{flushleft}\footnotesize The thirty acquisitions are repeated interrogations of nine sealed vials recorded within a single hour, so the effective sample size is the number of agents. On that unit the paired comparison between the two libraries gives, excluding HN2, $b = 4$ discordant pairs in favour of the measured library and $c = 0$ against, an exact two-sided $p = 0.12$ at the molecular level; including HN2, $b = 5$, $c = 0$, $p = 0.062$. At the class level, excluding HN2, the two libraries do not disagree on any agent. The molecular contrast is therefore a consistent direction on a sample too small to reach conventional significance, and we report it as such.\end{flushleft}
\end{table}
%%%%%%%%%%%%%%%%%%%%%%%%%%%%%%%%%%%%%%%%%%%%%%%%%%%%%%%%%%%%%%%%%%%%%
%%%%%%%%%%%%%%%%%%%%%%%%%%%%%%%%%%%%%%%%%%%%%%%%%%%%%%%%%%%%%%%%%%%%%
%%  table: hn2
%%%%%%%%%%%%%%%%%%%%%%%%%%%%%%%%%%%%%%%%%%%%%%%%%%%%%%%%%%%%%%%%%%%%%
\begin{table}[htbp]
\centering
\caption{Every reported quantity with and without the HN2 vial, whose contents are not established (Section~S5).}
\label{tab:si-hn2}
\small
\begin{tabular}{lrr}
\toprule
Quantity & All acquisitions & HN2 excluded \\
\midrule
Queries & 30 & 27 \\
Measured, molecular & 30 & 27 \\
Computed, molecular (10-member) & 13 & 13 \\
Computed, molecular (54-member) & 13 & 13 \\
Measured, class & 30 & 27 \\
Computed, class (10-member) & 25 & 25 \\
Computed, class (54-member) & 25 & 25 \\
Measured self-similarity & 0.983 & 0.981 \\
Computed self-similarity & 0.674 & 0.742 \\
Measured margin & 0.219 & 0.150 \\
Computed margin & 0.056 & 0.052 \\
\bottomrule
\end{tabular}

\begin{flushleft}\footnotesize Excluding HN2 raises the computed-library self-similarity from 0.674 to 0.742, so the deficit against the measured library falls from 0.309 to 0.239, and it removes the entire class-level difference between the two libraries. The HN2 acquisitions also carry an atypically large measured margin, which inflates the measured mean. We quote the HN2-excluded values in the abstract and the conclusions and retain the full set here.\end{flushleft}
\end{table}
%%%%%%%%%%%%%%%%%%%%%%%%%%%%%%%%%%%%%%%%%%%%%%%%%%%%%%%%%%%%%%%%%%%%%

\subsection{HN2}
\label{si:hn2}

The three HN2 acquisitions are assigned to a Novichok agent at a similarity to their own
computed entry of 0.057. A measured library identifies the HN2 vial correctly, because a
replicate of that vial is in it; only the computed library can state that the material is
not HN2. Those two results are simultaneously true only if the vial does not contain
HN2. No assay of the material is available and we do not claim to have established what
it contains; we report it as evidence about the sample rather than about the method.

\subsection{Where the computed library loses}

Figure~\ref{fig:si-spectra} shows the measured and calibrated computed spectrum of every
agent that has a computed entry, with the affected window shaded.
Table~\ref{tab:si-regions} resolves the measured-to-computed cosine by spectral region
against the measured-to-measured cosine in the same region, and splits the same
comparison by whether the agent contains phosphorus.
Table~\ref{tab:si-regbyagent} gives the per-agent values.

%%%%%%%%%%%%%%%%%%%%%%%%%%%%%%%%%%%%%%%%%%%%%%%%%%%%%%%%%%%%%%%%%%%%%
%%  figure: measured vs computed, all agents
%%%%%%%%%%%%%%%%%%%%%%%%%%%%%%%%%%%%%%%%%%%%%%%%%%%%%%%%%%%%%%%%%%%%%
\begin{figure}[htbp]
\centering
\includegraphics[width=\textwidth]{si_spectra.png}
\caption{Measured (upward, black) and calibrated computed (downward, red) Raman spectrum
of each of the eleven agents, all of which now have a computed library entry. The measured trace is the mean of that agent's genuine acquisitions and both
traces are normalized to unit maximum. The shaded band is the 600--850\,\wn{} window. The
number beside each agent is the cosine over the full 250--1800\,\wn{} range and the number
in the box is the cosine over the shaded window alone. The calibration is the one used
throughout, $s_{\mathrm{low}}$ = 1.016, $s_{\mathrm{high}}$ = 0.990 and $\Gamma$ = 22\,\wn{},
with no parameter fitted to identification. The window is where the six
phosphorus-containing agents lose and the three that contain none do not: compare
(b), (d) and (f) with (g), (h) and (i). HN2 (j) is the vial discussed in
Section~\ref{si:hn2} and is shown for completeness; it was never queried.}
\label{fig:si-spectra}
\end{figure}
%%%%%%%%%%%%%%%%%%%%%%%%%%%%%%%%%%%%%%%%%%%%%%%%%%%%%%%%%%%%%%%%%%%%%

%%%%%%%%%%%%%%%%%%%%%%%%%%%%%%%%%%%%%%%%%%%%%%%%%%%%%%%%%%%%%%%%%%%%%
%%  table: regions
%%%%%%%%%%%%%%%%%%%%%%%%%%%%%%%%%%%%%%%%%%%%%%%%%%%%%%%%%%%%%%%%%%%%%
\begin{table}[htbp]
\centering
\caption{Where the computed library loses, resolved by spectral region. The measured-to-computed cosine is compared with the measured-to-measured cosine in the same region, which is the ceiling that measurement repeatability permits. The last four columns split the same comparison by whether the agent contains phosphorus (seven agents) or not (three).}
\label{tab:si-regions}
\footnotesize
\resizebox{\textwidth}{!}{%
\begin{tabular}{lrrrrrrrr}
\toprule
Region (\wn) & Computed & Repeatability & Gap & Intensity & P-containing & No P & $\Delta$ & $p$ \\
\midrule
250--600 & 0.770 & 0.905 & 0.136 & 20.2\% & 0.810 & 0.675 & -0.135 & 0.967 \\
\textbf{600--850} & \textbf{0.719} & \textbf{0.982} & \textbf{0.263} & \textbf{29.2\%} & \textbf{0.636} & \textbf{0.912} & \textbf{+0.276} & \textbf{0.058} \\
850--1150 & 0.826 & 0.982 & 0.156 & 16.6\% & 0.840 & 0.793 & -0.047 & 0.808 \\
1150--1400 & 0.796 & 0.982 & 0.186 & 12.4\% & 0.814 & 0.755 & -0.058 & 0.867 \\
1400--1800 & 0.698 & 0.989 & 0.291 & 21.6\% & 0.758 & 0.557 & -0.201 & 1.000 \\
\bottomrule
\end{tabular}
}

\begin{flushleft}\footnotesize The largest gap coincides with the largest share of measured intensity. Within the 600--850\,\wn{} window the deficit belongs to the phosphorus-containing agents; in each of the four other windows the ordering reverses. With three agents on one side and seven on the other, the smallest attainable one-sided $p$ is 0.008 and the observed value is $0.058$, so the separation is not significant at the 0.05 level; the argument rests on the size of the effect and on its confinement to one window rather than on that $p$ value. $p$ is one-sided Mann--Whitney.\end{flushleft}
\end{table}
%%%%%%%%%%%%%%%%%%%%%%%%%%%%%%%%%%%%%%%%%%%%%%%%%%%%%%%%%%%%%%%%%%%%%
%%%%%%%%%%%%%%%%%%%%%%%%%%%%%%%%%%%%%%%%%%%%%%%%%%%%%%%%%%%%%%%%%%%%%
%%  table: regbyagent
%%%%%%%%%%%%%%%%%%%%%%%%%%%%%%%%%%%%%%%%%%%%%%%%%%%%%%%%%%%%%%%%%%%%%
\begin{table}[htbp]
\centering
\caption{Measured-to-computed cosine by region for every queried agent. Bold entries are the 600--850\,\wn{} window of the phosphorus-containing agents.}
\label{tab:si-regbyagent}
\footnotesize
\resizebox{\textwidth}{!}{%
\begin{tabular}{llrrrrr}
\toprule
Agent & Contains P & 250--600 skeletal/torsion & 600--850 & 850--1150 & 1150--1400 & 1400--1800 CH deform \\
\midrule
GD & yes & 0.882 & \textbf{0.884} & 0.883 & 0.943 & 0.759 \\
GF & yes & 0.596 & \textbf{0.323} & 0.928 & 0.921 & 0.622 \\
VX & yes & 0.775 & \textbf{0.964} & 0.845 & 0.737 & 0.789 \\
HD & no & 0.519 & 0.970 & 0.707 & 0.704 & 0.609 \\
HN1 & no & 0.784 & 0.910 & 0.768 & 0.715 & 0.533 \\
HN3 & no & 0.723 & 0.857 & 0.904 & 0.847 & 0.530 \\
A-230 & yes & 0.903 & \textbf{0.641} & 0.820 & 0.897 & 0.869 \\
A-232 & yes & 0.859 & \textbf{0.474} & 0.746 & 0.692 & 0.762 \\
A-234 & yes & 0.865 & \textbf{0.697} & 0.853 & 0.758 & 0.791 \\
A-242 & yes & 0.791 & \textbf{0.471} & 0.804 & 0.749 & 0.717 \\
\bottomrule
\end{tabular}
}

\begin{flushleft}\footnotesize The phosphorus-containing agents span 0.323 to 0.964 in the 600--850\,\wn{} window against 0.857--0.970 for the three that contain no phosphorus. Conformational flexibility does not order this: VX is the best of the phosphorus agents in that window and GF the worst, while a GFN2 conformer census over all eleven agents gives no association between conformational entropy and this cosine (Pearson $r$ = +0.23, $n$ = 10).\end{flushleft}
\end{table}
%%%%%%%%%%%%%%%%%%%%%%%%%%%%%%%%%%%%%%%%%%%%%%%%%%%%%%%%%%%%%%%%%%%%%
\clearpage

Table~\ref{tab:si-pfmult} makes the multiplicity of testing five windows explicit. No
window survives correction, and the smallest uncorrected $p$ belongs to a window in which
the ordering is reversed. We therefore report the phosphorus contrast as an effect size
with its sign structure and do not quote a $p$ value in support of it.

%%%%%%%%%%%%%%%%%%%%%%%%%%%%%%%%%%%%%%%%%%%%%%%%%%%%%%%%%%%%%%%%%%%%%
%%  table: pfmult
%%%%%%%%%%%%%%%%%%%%%%%%%%%%%%%%%%%%%%%%%%%%%%%%%%%%%%%%%%%%%%%%%%%%%
\begin{table}[htbp]
\centering
\caption{The phosphorus contrast in each of the five windows, with the multiplicity of testing five windows made explicit. Two-sided Mann--Whitney on seven phosphorus-containing against three phosphorus-free agents.}
\label{tab:si-pfmult}
\footnotesize
\resizebox{\textwidth}{!}{%
\begin{tabular}{lrrrrrr}
\toprule
Window (\wn) & P-containing & No P & $\Delta$ & $p$ & Bonferroni & Holm \\
\midrule
250-600 skeletal/torsion & 0.810 & 0.675 & -0.135 & 0.117 & 0.583 & 0.467 \\
600-850 & 0.636 & 0.912 & +0.276 & 0.117 & 0.583 & 0.467 \\
850-1150 & 0.840 & 0.793 & -0.047 & 0.517 & 1.000 & 0.767 \\
1150-1400 & 0.814 & 0.755 & -0.059 & 0.383 & 1.000 & 0.767 \\
1400-1800 CH deform & 0.758 & 0.557 & -0.201 & 0.017 & 0.083 & 0.083 \\
\bottomrule
\end{tabular}
}

\begin{flushleft}\footnotesize No window survives correction. The smallest uncorrected $p$ is not the 600--850\,\wn{} window but 1400--1800\,\wn{}, where the ordering is reversed. With six against three molecules the smallest two-sided $p$ the design can produce is 0.024, so this test has essentially no power and we do not quote a $p$ value in support of the localisation. The argument rests on the effect size, $\Delta = +0.276$ in the 600--850\,\wn{} window against $|\Delta| \leq 0.201$ elsewhere, and on the sign reversal in all four other windows.\end{flushleft}
\end{table}
%%%%%%%%%%%%%%%%%%%%%%%%%%%%%%%%%%%%%%%%%%%%%%%%%%%%%%%%%%%%%%%%%%%%%
\clearpage

\subsection{The Novichok agents}

Table~\ref{tab:si-nov} compares the four A-series agents with one another and with
themselves in each library.

%%%%%%%%%%%%%%%%%%%%%%%%%%%%%%%%%%%%%%%%%%%%%%%%%%%%%%%%%%%%%%%%%%%%%
%%  table: nov
%%%%%%%%%%%%%%%%%%%%%%%%%%%%%%%%%%%%%%%%%%%%%%%%%%%%%%%%%%%%%%%%%%%%%
\begin{table}[htbp]
\centering
\caption{The four Novichok agents compared with one another and with themselves, in each library. \emph{Sibling} is the highest similarity to another A-series agent.}
\label{tab:si-nov}
\small
\begin{tabular}{lrrrr}
\toprule
Agent & Measured self & Measured sibling & Computed self & Computed sibling \\
\midrule
A-230 & 0.970 & 0.847 & 0.801 & 0.928 \\
A-232 & 1.000 & 0.881 & 0.652 & 0.928 \\
A-234 & 0.941 & 0.881 & 0.722 & 0.925 \\
A-242 & 0.958 & 0.710 & 0.706 & 0.844 \\
\bottomrule
\end{tabular}

\begin{flushleft}\footnotesize The mutual similarities are comparable in the two libraries (0.71--0.88 measured against 0.84--0.93 computed), while the similarity to the correct molecule is not. The four agents share the P(=O)(F)--N=C amidine core and differ only in the phosphorus substituent and, for A-242, in guanidine substitution. The computed library therefore fails on them not because the congeners are too much alike but because the correct entry is not drawn well enough.\end{flushleft}
\end{table}
%%%%%%%%%%%%%%%%%%%%%%%%%%%%%%%%%%%%%%%%%%%%%%%%%%%%%%%%%%%%%%%%%%%%%
\clearpage

%%%%%%%%%%%%%%%%%%%%%%%%%%%%%%%%%%%%%%%%%%%%%%%%%%%%%%%%%%%%%%%%%%%%%
\begin{table}[htbp]
\centering
\caption{Beyond top-1: where the correct entry actually ranks. A library search returns a
hit list, so the rank of the correct entry is the quantity a practitioner reads.}
\label{tab:si-topk}
\small
\begin{tabular}{lrr}
\toprule
 & Computed library (54) & Measured library (11) \\
\midrule
Top-1 & 13/30 & 30/30 \\
Top-3 & 16/30 & 30/30 \\
Top-5 & 22/30 & 30/30 \\
Top-10 & 22/30 & 30/30 \\
\midrule
Median rank of correct entry & 2 & 1 \\
Worst rank of correct entry & 46 & 1 \\
\bottomrule
\end{tabular}

\begin{flushleft}\footnotesize Rank distribution for the computed library: rank 1 for 13
acquisitions, rank 2 for 3, rank 4 for 3, rank 5 for 3, rank 11 for 1, rank 14 for 4, and
ranks 43, 45 and 46 for one each. Allowing three candidates instead of one recovers three
of the seventeen errors and no more; for seven of the thirty the correct entry is not in
the top ten. The deficit is therefore not a matter of narrow misses throughout.

It is a matter of narrow misses in part. Six of the seventeen errors are decided by a
margin below 0.01: the three A-234 acquisitions assigned to A-232 at margins of 0.0003 to
0.0004, with the correct entry at rank 4, and the three HN1 acquisitions assigned to HN3
at margins of 0.0039 to 0.0042, with the correct entry at rank 2. Ties at this precision
are not meaningfully resolved by the metric, and we report them as such rather than as
six independent failures. The remaining eleven errors have the correct entry at rank 4 or
worse with margins above 0.01.\end{flushleft}
\end{table}

\section{Is the deficit an artifact of the metric or of the lineshape?}
\label{si:metric}

Two explanations of the result do not involve the calculation at all. The first is that
full-profile cosine is the wrong instrument: at a broadening of 22\,\wn{} a cosine is close
to a monotone function of the band-position residual, and the residual reported here is 0.6
of that width, where the sensitivity is greatest. The second is that the deficit is a
lineshape mismatch, since a single global $\Gamma$ is applied to library entries whose
measured widths span 9.2 to 48.4\,\wn{}. Both are tested here and neither survives.

\subsection{Twelve metrics, with the measured library as a control}

Table~\ref{tab:si-metrics} scores the same thirty queries against the same libraries with
twelve similarity metrics. Measured and computed spectra pass through the identical peak
detector on the identical grid, and the measured library is carried through the same code,
so that a metric that performs poorly on the computed library can be distinguished from a
metric that is badly implemented.

%%%%%%%%%%%%%%%%%%%%%%%%%%%%%%%%%%%%%%%%%%%%%%%%%%%%%%%%%%%%%%%%%%%%%
%%  table: metrics
%%%%%%%%%%%%%%%%%%%%%%%%%%%%%%%%%%%%%%%%%%%%%%%%%%%%%%%%%%%%%%%%%%%%%
\begin{table}[htbp]
\centering
\caption{The same thirty queries scored with twelve different similarity metrics. The measured library is carried through the identical code as a control, so that a metric that performs poorly on the computed library can be distinguished from a metric that is simply badly implemented.}
\label{tab:si-metrics}
\footnotesize
\resizebox{\textwidth}{!}{%
\begin{tabular}{lcccccc}
\toprule
Metric & Meas.\ top-1 & Meas.\ self & Comp.\ top-1 (11) & Comp.\ top-1 (54) & Comp.\ by agent & Comp.\ self \\
\midrule
Full-profile cosine, $\Gamma$ = 22\,\wn{} (adopted) & 30/30 & 0.983 & 13/30 & 13/30 & 4/9 & 0.674 \\
Full-profile cosine, $\Gamma$ = 30\,\wn{} & n/a & n/a & 16/30 & 16/30 & 5/9 & 0.688 \\
Full-profile cosine, $\Gamma$ = 40\,\wn{} & n/a & n/a & 16/30 & 16/30 & 5/9 & 0.692 \\
Full-profile cosine, $\Gamma$ = 50\,\wn{} & n/a & n/a & 16/30 & 16/30 & 5/9 & 0.688 \\
Full-profile cosine, per-agent measured $\Gamma$ & n/a & n/a & 12/30 & 12/30 & 4/9 & 0.670 \\
First-derivative cosine & 30/30 & 0.981 & 10/30 & 10/30 & 3/9 & 0.350 \\
Shift-tolerant cosine, $\pm$3 pixels ($\pm$17\,\wn{}) & 30/30 & 0.983 & 16/30 & 16/30 & 5/9 & 0.720 \\
Peak-list hit quality, 10\,\wn{} tolerance & 28/30 & 0.801 & 8/30 & 4/30 & 2/9 & 0.122 \\
Peak-list hit quality, 20\,\wn{} tolerance & 29/30 & 0.835 & 11/30 & 11/30 & 3/9 & 0.232 \\
Peak-list hit quality, 35\,\wn{} tolerance & 29/30 & 0.856 & 11/30 & 11/30 & 3/9 & 0.326 \\
Band-position Dice, 20\,\wn{} tolerance & 28/30 & 0.891 & 8/30 & 6/30 & 2/9 & 0.459 \\
Band-position Dice, 35\,\wn{} tolerance & 25/30 & 0.898 & 10/30 & 9/30 & 2/9 & 0.622 \\
\bottomrule
\end{tabular}
}

\begin{flushleft}\footnotesize No metric rescues the computed library. The adopted metric gives 13 of 30, the best of the twelve is 16 of 30 (Full-profile cosine), and every peak-list metric, which is what commercial library search actually uses, lies between 8 and 11 of 30. The measured library reaches 25 to 30 of 30 under all twelve, so the peak-list implementations are sound and the collapse is a property of the computed library and not of the comparison. The similarity metric is therefore excluded as an explanation of the result.\end{flushleft}
\end{table}
%%%%%%%%%%%%%%%%%%%%%%%%%%%%%%%%%%%%%%%%%%%%%%%%%%%%%%%%%%%%%%%%%%%%%
\clearpage

\subsection{Which attribute of the computed line list carries the deficit}

Table~\ref{tab:si-decomp} replaces one attribute of the computed line list at a time with
the measured value and resynthesizes the spectrum in the 600--850\,\wn{} window, taking the
measured spectrum itself as the reference. Table~\ref{tab:si-unmatched} then quantifies
what the substitutions cannot reach.

%%%%%%%%%%%%%%%%%%%%%%%%%%%%%%%%%%%%%%%%%%%%%%%%%%%%%%%%%%%%%%%%%%%%%
%%  table: decomp
%%%%%%%%%%%%%%%%%%%%%%%%%%%%%%%%%%%%%%%%%%%%%%%%%%%%%%%%%%%%%%%%%%%%%
\begin{table}[htbp]
\centering
\caption{Decomposition of the 600--850\,\wn{} deficit. One attribute of the computed line list at a time is replaced by the measured value and the spectrum resynthesized; the reference is the measured spectrum itself.}
\label{tab:si-decomp}
\footnotesize
\resizebox{\textwidth}{!}{%
\begin{tabular}{llcrrrrrr}
\toprule
Agent & P & Paired/comp./meas. & Baseline & Positions & Intensities & Linewidth & Pos.+int. & All three \\
\midrule
GD & $\bullet$ & 6/15/6 & 0.884 & 0.872 & 0.939 & 0.875 & 0.956 & 0.965 \\
GF & $\bullet$ & 4/10/5 & 0.314 & 0.327 & 0.397 & 0.274 & 0.421 & 0.387 \\
HN1 & $\circ$ & 3/9/3 & 0.910 & 0.960 & 0.926 & 0.909 & 0.980 & 0.980 \\
HN2 & $\circ$ & 0/3/1 & 0.394 & 0.394 & 0.394 & 0.345 & 0.394 & 0.345 \\
HN3 & $\circ$ & 3/7/3 & 0.858 & 0.898 & 0.949 & 0.858 & 0.980 & 0.979 \\
A-230 & $\bullet$ & 10/13/12 & 0.638 & 0.938 & 0.633 & 0.574 & 0.936 & 0.951 \\
A-232 & $\bullet$ & 5/7/7 & 0.471 & 0.908 & 0.486 & 0.493 & 0.974 & 0.967 \\
A-234 & $\bullet$ & 3/3/4 & 0.703 & 0.815 & 0.828 & 0.734 & 0.977 & 0.977 \\
A-242 & $\bullet$ & 9/13/11 & 0.446 & 0.513 & 0.478 & 0.537 & 0.585 & 0.598 \\
\textbf{All} &  &  & \textbf{0.624} & \textbf{0.736} & \textbf{0.670} & \textbf{0.622} & \textbf{0.800} & \textbf{0.794} \\
\textbf{P-containing} &  &  & \textbf{0.576} & \textbf{0.729} & \textbf{0.627} & \textbf{0.581} & \textbf{0.808} & \textbf{0.807} \\
\textbf{No P} &  &  & \textbf{0.720} & \textbf{0.750} & \textbf{0.757} & \textbf{0.704} & \textbf{0.785} & \textbf{0.768} \\
\bottomrule
\end{tabular}
}

\begin{flushleft}\footnotesize The baseline column is rebuilt from the lines that fall in and around the window rather than taken from the full-spectrum synthesis, so it differs from Table~\ref{tab:si-regbyagent} by up to 0.03; the comparison of interest is the change within each row. Replacing the linewidth with the width measured for that agent changes the window cosine by -0.002 on average, so the single global $\Gamma$ is not the origin of the deficit. Position contributes about twice what intensity does among the phosphorus-containing agents (+0.153 against +0.051). Replacing all three attributes leaves the cosine at 0.794 rather than at unity, so the largest remaining component is neither position nor intensity nor width but the set of bands itself (Table~\ref{tab:si-unmatched}).\end{flushleft}
\end{table}
%%%%%%%%%%%%%%%%%%%%%%%%%%%%%%%%%%%%%%%%%%%%%%%%%%%%%%%%%%%%%%%%%%%%%
%%%%%%%%%%%%%%%%%%%%%%%%%%%%%%%%%%%%%%%%%%%%%%%%%%%%%%%%%%%%%%%%%%%%%
%%  table: unmatched
%%%%%%%%%%%%%%%%%%%%%%%%%%%%%%%%%%%%%%%%%%%%%%%%%%%%%%%%%%%%%%%%%%%%%
\begin{table}[htbp]
\centering
\caption{Bands inside the 600--850\,\wn{} window that find no partner in the other library within 35\,\wn{}, expressed as a share of the intensity carried in that window.}
\label{tab:si-unmatched}
\small
\resizebox{\textwidth}{!}{%
\begin{tabular}{lrrrrr}
\toprule
Agent & Computed bands & Measured bands & Paired & Computed with no measured partner & Measured with no computed partner \\
\midrule
GD & 4 & 3 & 3 & 10\% & 0\% \\
GF & 4 & 4 & 3 & 50\% & 6\% \\
HN1 & 5 & 3 & 3 & 10\% & 0\% \\
HN2 & 2 & 0 & 0 & 100\% & 100\% \\
HN3 & 5 & 2 & 2 & 32\% & 0\% \\
A-230 & 7 & 6 & 6 & 8\% & 0\% \\
A-232 & 2 & 2 & 1 & 75\% & 88\% \\
A-234 & 2 & 2 & 2 & 0\% & 0\% \\
A-242 & 7 & 5 & 4 & 31\% & 38\% \\
\textbf{Mean} & \textbf{4.2} & \textbf{3.0} & \textbf{2.7} & \textbf{35\%} & \textbf{26\%} \\
\bottomrule
\end{tabular}
}

\begin{flushleft}\footnotesize Measured bands almost always find a computed partner: the share is zero for GD, HN1, HN3, A-230 and A-234. The asymmetry runs the other way. The calculation places on average 35\% of the window intensity in bands the measurement does not show, rising to 50\% for GF, which is one of the misassigned agents. The deficit in this window is therefore not that the calculation misplaces the measured bands but that it distributes intensity over a band set the measurement does not have. A-232 and HN2 are anomalous in both directions and are discussed separately in Sections~S5 and~S8.\end{flushleft}
\end{table}
%%%%%%%%%%%%%%%%%%%%%%%%%%%%%%%%%%%%%%%%%%%%%%%%%%%%%%%%%%%%%%%%%%%%%
\clearpage

%%%%%%%%%%%%%%%%%%%%%%%%%%%%%%%%%%%%%%%%%%%%%%%%%%%%%%%%%%%%%%%%%%%%%
\begin{table}[htbp]
\centering
\caption{The preprocessing applied symmetrically. The baseline correction and the
dead-pixel interpolation are applied to the measurement only in the adopted pipeline,
because the computed spectrum is synthesised and has no baseline. A Lorentzian line list
does, however, build a pseudo-continuum from its own wings, which the measured spectrum
has had removed. Here the identical AsLS ($\lambda = 10^{4}$, $p = 0.01$) and the same
three interpolated pixels are applied to the synthesised computed spectra as well.}
\label{tab:si-aslssym}
\small
\begin{tabular}{lrrrrrcr}
\toprule
Computed spectra & 250--600 & 600--850 & 850--1150 & 1150--1400 & 1400--1800 & Molecular & Self cos. \\
\midrule
As adopted (no AsLS) & 0.759 & 0.690 & 0.761 & 0.742 & 0.647 & 13/30 & 0.674 \\
AsLS applied & 0.700 & 0.657 & 0.729 & 0.678 & 0.612 & 13/30 & 0.644 \\
AsLS + dead pixels & 0.700 & 0.657 & 0.729 & 0.678 & 0.612 & 13/30 & 0.644 \\
\midrule
\textbf{Change} & $-$0.059 & $-$0.034 & $-$0.033 & $-$0.064 & $-$0.035 & \textbf{0} & $-$0.030 \\
\bottomrule
\end{tabular}

\begin{flushleft}\footnotesize Removing the Lorentzian pseudo-continuum from the computed
spectra lowers the agreement in every window and leaves the identification at 13 of 30.
The asymmetry in the adopted pipeline therefore works slightly in the computed library's
favour, not against it, and it does not generate the deficit this paper reports. The
dead-pixel treatment changes nothing at three decimal places, which is expected: the
three pixels lie at 813.7, 1220.6 and 1634.4\,\wn{} and none is near a dominant computed
band.\end{flushleft}
\end{table}

\section{Abstention and the class-absent condition}
\label{si:absten}

\subsection{Deleting one molecule}

Table~\ref{tab:si-loo} deletes every reference spectrum of the query's own molecule from
the measured library and re-queries. The table reports the number of library members before
and after the deletion, so that the ablation is verifiable from the table itself; no query
is assigned to its own molecule. Every query except HN2 falls to a member of its own class,
at an average similarity of 0.83 against the within-molecule value of 0.98: with the
congeners still present, the assignment degrades gracefully. HN2 falls to a Novichok agent,
consistent with Section~\ref{si:hn2}.

%%%%%%%%%%%%%%%%%%%%%%%%%%%%%%%%%%%%%%%%%%%%%%%%%%%%%%%%%%%%%%%%%%%%%
%%  table: loo
%%%%%%%%%%%%%%%%%%%%%%%%%%%%%%%%%%%%%%%%%%%%%%%%%%%%%%%%%%%%%%%%%%%%%
\begin{table}[htbp]
\centering
\caption{Leave-one-molecule-out against the measured library. Each agent's own reference spectra were deleted before matching, so the query must fall to a different molecule.}
\label{tab:si-loo}
\small
\begin{tabular}{llrp{3.6cm}rrr}
\toprule
Agent & $n$ & Library & Falls to & Class kept & Top-1 cosine & Margin \\
\midrule
GD & 4 & 11 $\rightarrow$ 10 & VX ($\times$4) & 4/4 & 0.789 & 0.079 \\
GF & 3 & 11 $\rightarrow$ 10 & GD ($\times$3) & 3/3 & 0.710 & 0.063 \\
HN1 & 3 & 11 $\rightarrow$ 10 & HN3 ($\times$3) & 3/3 & 0.935 & 0.222 \\
HN2 & 3 & 11 $\rightarrow$ 10 & A-242 ($\times$3) & 0/3 & 0.168 & 0.052 \\
HN3 & 3 & 11 $\rightarrow$ 10 & HN1 ($\times$3) & 3/3 & 0.934 & 0.197 \\
A-230 & 3 & 11 $\rightarrow$ 10 & A-234 ($\times$3) & 3/3 & 0.837 & 0.039 \\
A-232 & 3 & 11 $\rightarrow$ 10 & A-234 ($\times$3) & 3/3 & 0.881 & 0.074 \\
A-234 & 5 & 11 $\rightarrow$ 10 & A-232 ($\times$4), A-230 & 5/5 & 0.857 & 0.048 \\
A-242 & 3 & 11 $\rightarrow$ 10 & A-234 ($\times$3) & 3/3 & 0.697 & 0.031 \\
\textbf{All} & \textbf{30} &  &  & \textbf{27/30} & \textbf{0.764} & \textbf{0.087} \\
\textbf{All, HN2 excluded} & \textbf{27} &  &  & \textbf{27/27} & \textbf{0.830} & \textbf{0.090} \\
\bottomrule
\end{tabular}

\begin{flushleft}\footnotesize The \emph{Library} column reports the number of library members before and after the deletion, so that the ablation is verifiable from the table itself; no query is assigned to its own molecule. Every query except HN2 falls to a member of its own toxicological class, and the similarity to that neighbor (0.83 on average) is far below the within-molecule value of 0.98. This is the measured-library counterpart of the class-absent test: deleting one molecule leaves its congeners and the assignment degrades gracefully, whereas deleting the entire class does not (Table~\ref{tab:si-loco}). HN2 falls to a Novichok agent, consistent with the anomaly reported in Section~S5.\end{flushleft}
\end{table}
%%%%%%%%%%%%%%%%%%%%%%%%%%%%%%%%%%%%%%%%%%%%%%%%%%%%%%%%%%%%%%%%%%%%%
\clearpage

\subsection{Deleting the whole class}

Table~\ref{tab:si-loco} deletes every library member of the query's toxicological class,
so the correct answer is absent by construction and the workflow should decline. It does
not decline. It returns confident assignments whose similarities and margins fall inside
the ranges it returns when the answer is present.

%%%%%%%%%%%%%%%%%%%%%%%%%%%%%%%%%%%%%%%%%%%%%%%%%%%%%%%%%%%%%%%%%%%%%
%%  table: loco
%%%%%%%%%%%%%%%%%%%%%%%%%%%%%%%%%%%%%%%%%%%%%%%%%%%%%%%%%%%%%%%%%%%%%
\begin{table}[htbp]
\centering
\caption{Leave-one-class-out. Every library member of the query's toxicological class was deleted before matching, so the correct answer is absent by construction and the workflow should decline.}
\label{tab:si-loco}
\small
\begin{tabular}{llp{4.4cm}rrr}
\toprule
Agent & $n$ & Assigned to & Top-1 cosine & Margin & Library $z$ \\
\midrule
GD & 4 & EXP\_HN2 ($\times$4) & 0.755 & 0.003 & +1.27 \\
GF & 3 & EXP\_HN2 ($\times$3) & 0.588 & 0.024 & -0.17 \\
HN1 & 3 & EtGB ($\times$2), GB\_sarin & 0.699 & 0.000 & -0.01 \\
HN2 & 3 & EXP\_A242 ($\times$3) & 0.207 & 0.043 & -7.16 \\
HN3 & 3 & GB\_sarin ($\times$3) & 0.646 & 0.016 & -0.54 \\
A-230 & 3 & OP\_nPr\_iBu ($\times$3) & 0.757 & 0.030 & -0.11 \\
A-232 & 3 & OP\_nPr\_iBu ($\times$3) & 0.681 & 0.021 & -1.10 \\
A-234 & 5 & OP\_Et\_iBu ($\times$3), OP\_Me\_Me ($\times$2) & 0.767 & 0.027 & +0.40 \\
A-242 & 3 & OP\_sBu\_Me ($\times$3) & 0.590 & 0.005 & -1.87 \\
\bottomrule
\end{tabular}

\begin{flushleft}\footnotesize Class-level accuracy under this condition is 0 by construction. The workflow does not signal the absence: the similarities and margins it returns fall inside the ranges it returns when the answer is present, which is why no threshold on them separates the two conditions (Table~\ref{tab:si-absten}).\end{flushleft}
\end{table}
%%%%%%%%%%%%%%%%%%%%%%%%%%%%%%%%%%%%%%%%%%%%%%%%%%%%%%%%%%%%%%%%%%%%%
\clearpage

\subsection{Whether any score can tell the two apart}

Table~\ref{tab:si-absten} evaluates six scores under both conditions and against both
libraries, with thresholds set by nested leave-one-agent-out.

%%%%%%%%%%%%%%%%%%%%%%%%%%%%%%%%%%%%%%%%%%%%%%%%%%%%%%%%%%%%%%%%%%%%%
%%  table: absten
%%%%%%%%%%%%%%%%%%%%%%%%%%%%%%%%%%%%%%%%%%%%%%%%%%%%%%%%%%%%%%%%%%%%%
\begin{table}[htbp]
\centering
\caption{Six abstention scores evaluated on the same queries under three conditions: the class present in the library (accept), the query's entire toxicological class removed (decline), and the HN2 acquisitions. Thresholds are set by nested leave-one-agent-out, so no threshold is chosen on the queries it is applied to.}
\label{tab:si-absten}
\footnotesize
\resizebox{\textwidth}{!}{%
\begin{tabular}{lrrrrcrcrr}
\toprule
Score & In-library & Class absent & HN2 & AUC & Separable? & AUC (meas.) & Sep.? & Decline IN & Decline LOCO \\
\midrule
top-1 cosine & 0.780 & 0.699 & 0.206 & 0.829 & no & 1.000 & \textbf{yes} & 48\% & 93\% \\
margin & 0.038 & 0.018 & 0.043 & 0.775 & no & 0.909 & no & 59\% & 96\% \\
peak-to-sidelobe ratio & 2.843 & 2.095 & 4.657 & 0.852 & no & 0.985 & no & 30\% & 96\% \\
library-normalized $z$ & 0.782 & -0.188 & -7.165 & 0.775 & no & 1.000 & \textbf{yes} & 56\% & 93\% \\
unmixing residual & 0.553 & 0.681 & 0.978 & 0.798 & no & 1.000 & \textbf{yes} & 44\% & 93\% \\
band-match fraction & 0.833 & 0.538 & 0.333 & 0.770 & no & 0.824 & no & 74\% & 93\% \\
\bottomrule
\end{tabular}
}

\begin{flushleft}\footnotesize Against the computed library no score separates the two conditions. Driving the class-absent acceptance to zero requires a conjunction of the library-normalized $z$ and the top-1 cosine, and costs 63\% of the legitimate queries. Against the measured library three scores separate completely and a top-1 cosine threshold of 0.729 accepts 27 of 27 and declines 27 of 27. That threshold is an upper bound obtained on one instrument, one session and one vial per agent, and is not a transferable operating value; the contrast between the libraries is the result.\end{flushleft}
\end{table}
%%%%%%%%%%%%%%%%%%%%%%%%%%%%%%%%%%%%%%%%%%%%%%%%%%%%%%%%%%%%%%%%%%%%%
\clearpage

Because the thirty acquisitions interrogate nine vials, an area under the curve computed
over acquisitions overstates the evidence. Table~\ref{tab:si-aucci} recomputes it with the
agent as the resampling unit and attaches bootstrap intervals.

%%%%%%%%%%%%%%%%%%%%%%%%%%%%%%%%%%%%%%%%%%%%%%%%%%%%%%%%%%%%%%%%%%%%%
%%  table: aucci
%%%%%%%%%%%%%%%%%%%%%%%%%%%%%%%%%%%%%%%%%%%%%%%%%%%%%%%%%%%%%%%%%%%%%
\begin{table}[htbp]
\centering
\caption{Separation between the class-present and the class-absent condition, with the agent as the resampling unit. Intervals are 5000 bootstrap replicates over agents.}
\label{tab:si-aucci}
\footnotesize
\resizebox{\textwidth}{!}{%
\begin{tabular}{lcccc}
\toprule
Score & Measured AUC & 95\% CI & Computed AUC & 95\% CI \\
\midrule
top-1 cosine & 1.000 & [1.000, 1.000] & 0.766 & [0.639, 0.927] \\
margin & 0.923 & [0.829, 1.000] & 0.763 & [0.547, 0.976] \\
peak-to-sidelobe & 0.900 & [0.782, 1.000] & 0.790 & [0.596, 1.000] \\
library-normalized $z$ & 1.000 & [1.000, 1.000] & 0.723 & [0.618, 0.886] \\
unmixing residual & 1.000 & [1.000, 1.000] & 0.741 & [0.619, 0.904] \\
band-match fraction & 0.848 & [0.741, 0.957] & 0.729 & [0.542, 0.895] \\
\bottomrule
\end{tabular}
}

\begin{flushleft}\footnotesize Resampling agents rather than acquisitions removes the threefold inflation that treating the thirty acquisitions as independent would introduce. On the measured library three scores separate the two conditions completely and the bootstrap never breaks that separation. On the computed library the best score reaches 0.790 and no score separates the two conditions on the observed data; the intervals are wide, so this is a statement about what nine agents can establish and not a proof that no score could separate them.\end{flushleft}
\end{table}
%%%%%%%%%%%%%%%%%%%%%%%%%%%%%%%%%%%%%%%%%%%%%%%%%%%%%%%%%%%%%%%%%%%%%
\clearpage

%%%%%%%%%%%%%%%%%%%%%%%%%%%%%%%%%%%%%%%%%%%%%%%%%%%%%%%%%%%%%%%%%%%%%
\begin{table}[htbp]
\centering
\caption{Abstention with the candidate count equalised. In the adopted analysis the
computed library offers 54 candidates and the measured library 11, and the top-1 score is
an extreme-value statistic over the candidate set. Here the computed library is restricted
to the same eleven measured agents so that both conditions present the same number of
candidates.}
\label{tab:si-abstsize}
\small
\begin{tabular}{lrrcc}
\toprule
 & \multicolumn{2}{c}{AUC, in-library vs class-absent} & \multicolumn{2}{c}{Separates completely?} \\
\cmidrule(lr){2-3}\cmidrule(lr){4-5}
Score & Computed (11) & Measured (11) & Computed & Measured \\
\midrule
Top-1 cosine & 0.860 & 1.000 & no & \textbf{yes} \\
Margin & 0.689 & 0.909 & no & no \\
Peak-to-sidelobe ratio & 0.554 & 0.985 & no & no \\
Library-normalized $z$ & 0.746 & 1.000 & no & \textbf{yes} \\
Unmixing residual & 0.844 & 1.000 & no & \textbf{yes} \\
Band-match fraction & 0.746 & 0.824 & no & no \\
\bottomrule
\end{tabular}

\begin{flushleft}\footnotesize With the candidate counts equal, no score separates the two
conditions against the computed library and three separate them completely against the
measured library, which is the same conclusion the 54-candidate analysis reaches. The
single ordinary top-1 cosine threshold of 0.729 still accepts 27 of 27 in-library queries
and declines 27 of 27 class-absent ones against the measured library. The candidate count
is therefore not what produces the contrast. Two scores move appreciably under the
restriction --- the peak-to-sidelobe ratio loses most of its AUC, from 0.852 to 0.554,
because it is defined against the spread of the remaining candidates --- but no score
crosses the separability threshold in either direction.\end{flushleft}
\end{table}

\section{Raising the level of theory}
\label{si:theory}

Three corrections were tested against the diagnosis of Section~\ref{si:ident}: implicit
solvation with CPCM (acetone, chosen for a dielectric constant near that of the neat
agents), promotion of phosphorus, fluorine, sulfur and chlorine to def2-TZVPP with the
remaining atoms at 6-311++G(d,p), and the range-separated functional $\omega$B97X-D3.
Geometries were re-optimized in every case. Each improved calculation was swapped into
the library and the frequency scale refitted independently, so that no improvement is
understated by an inherited calibration. The ten trials are given in
Table~\ref{tab:si-theory} and the calculations themselves in
Table~\ref{tab:si-b2jobs}.

%%%%%%%%%%%%%%%%%%%%%%%%%%%%%%%%%%%%%%%%%%%%%%%%%%%%%%%%%%%%%%%%%%%%%
%%  table: theory
%%%%%%%%%%%%%%%%%%%%%%%%%%%%%%%%%%%%%%%%%%%%%%%%%%%%%%%%%%%%%%%%%%%%%
\begin{table}[htbp]
\centering
\caption{Ten agent--correction trials. Each improved calculation was swapped into the library and the frequency scale refitted independently, so that no improvement is understated by an inherited calibration. All ten trials were computed from the pre-correction geometries, and the baseline column is therefore the pre-correction library; the geometry correction reported in Section~S9 is a separate effect and the two are not additive.}
\label{tab:si-theory}
\footnotesize
\resizebox{\textwidth}{!}{%
\begin{tabular}{llrrrrrc}
\toprule
Agent & Correction & 600--850 base & Improved & $\Delta$ & Self base & Self new & Top-1 \\
\midrule
GF & $\omega$B97X-D3 & 0.548 & \textbf{0.918} & \textbf{+0.370} & 0.643 & 0.793 & $\times$ $\rightarrow$ \checkmark \\
GF & tight P/F/S/Cl & 0.548 & \textbf{0.832} & \textbf{+0.284} & 0.643 & 0.797 & $\times$ $\rightarrow$ \checkmark \\
A-232 & $\omega$B97X-D3 & 0.199 & \textbf{0.315} & \textbf{+0.116} & 0.606 & 0.564 & $\times$ $\rightarrow$ --- \\
A-230 & tight P/F/S/Cl & 0.627 & 0.694 & +0.068 & 0.810 & 0.811 & \checkmark $\rightarrow$ \checkmark \\
A-232 & CPCM & 0.199 & 0.239 & +0.040 & 0.606 & 0.659 & $\times$ $\rightarrow$ --- \\
A-232 & tight P/F/S/Cl & 0.199 & 0.226 & +0.027 & 0.606 & 0.616 & $\times$ $\rightarrow$ --- \\
GF & CPCM & 0.548 & 0.552 & +0.005 & 0.643 & 0.623 & $\times$ $\rightarrow$ --- \\
HD & tight P/F/S/Cl & 0.969 & 0.959 & -0.010 & 0.868 & 0.865 & \checkmark $\rightarrow$ \checkmark \\
GD & tight P/F/S/Cl & 0.888 & \textbf{0.726} & \textbf{-0.161} & 0.859 & 0.766 & \checkmark $\rightarrow$ \checkmark \\
A-234 & tight P/F/S/Cl & 0.673 & \textbf{0.485} & \textbf{-0.188} & 0.759 & 0.724 & $\times$ $\rightarrow$ --- \\
\bottomrule
\end{tabular}
}

\begin{flushleft}\footnotesize Of the 10 pairs, 3 improved the window substantially ($\Delta > 0.10$), 2 improved modestly, 3 were unchanged and 2 degraded. Across all 10 the refitted scale factors stay within 1.012--1.026 and 0.988--0.992, and the global band-position RMSE between 10.2 and 11.4\,\wn{}, so no global adjustment reaches these errors.\end{flushleft}
\end{table}
%%%%%%%%%%%%%%%%%%%%%%%%%%%%%%%%%%%%%%%%%%%%%%%%%%%%%%%%%%%%%%%%%%%%%
\clearpage

%%%%%%%%%%%%%%%%%%%%%%%%%%%%%%%%%%%%%%%%%%%%%%%%%%%%%%%%%%%%%%%%%%%%%
%%  table: b2 jobs
%%%%%%%%%%%%%%%%%%%%%%%%%%%%%%%%%%%%%%%%%%%%%%%%%%%%%%%%%%%%%%%%%%%%%
\begin{table}[htbp]
\centering
\caption{The ten higher-level calculations. Every job converged to a true minimum and terminated normally, and the number of real modes agrees with the corresponding baseline calculation in Table~\ref{tab:si-freq} in every case. Wall time is for eight MPI processes.}
\label{tab:si-b2jobs}
\footnotesize
\begin{tabular}{llrrrr}
\toprule
Agent & Correction & Atoms & $N_\mathrm{real}$ & $N_\mathrm{imag}$ & Wall time (h) \\
\midrule
GF & CPCM (acetone, $\varepsilon$ = 21.06) & 25 & 69 & 0 & 6.4 \\
GF & tight P/F/S/Cl & 25 & 69 & 0 & 6.7 \\
GF & $\omega$B97X-D3 & 25 & 69 & 0 & 10.5 \\
GD & tight P/F/S/Cl & 27 & 75 & 0 & 6.4 \\
HD & tight P/F/S/Cl & 15 & 39 & 0 & 1.9 \\
A-230 & tight P/F/S/Cl & 28 & 78 & 0 & 10.9 \\
A-232 & CPCM (acetone, $\varepsilon$ = 21.06) & 29 & 81 & 0 & 13.4 \\
A-232 & tight P/F/S/Cl & 29 & 81 & 0 & 9.7 \\
A-232 & $\omega$B97X-D3 & 29 & 81 & 0 & 18.3 \\
A-234 & tight P/F/S/Cl & 32 & 90 & 0 & 11.8 \\
\bottomrule
\end{tabular}

\begin{flushleft}\footnotesize All ten were run with ORCA 6.1.1, \texttt{TightOpt} followed by \texttt{NumFreq} with \texttt{\%elprop Polar 1}, on eight MPI processes with 1400 MB per core. The tight-basis runs promote phosphorus, fluorine, sulfur and chlorine to def2-TZVPP through \texttt{\%basis NewGTO} and leave every other atom at 6-311++G(d,p); the CPCM runs add \texttt{CPCM(Acetone)} to the baseline keywords; the $\omega$B97X-D3 runs replace the functional and change nothing else. CPCM was applied to the two agents with the lowest window cosine as a negative control for the condensed phase, and the range-separated functional to the same two agents.\end{flushleft}
\end{table}
%%%%%%%%%%%%%%%%%%%%%%%%%%%%%%%%%%%%%%%%%%%%%%%%%%%%%%%%%%%%%%%%%%%%%
\clearpage

Table~\ref{tab:si-bandshift} gives the mechanism. The basis-set change moves the dominant
600--850\,\wn{} band upward in five of the six molecules and leaves it essentially where it
was in the sixth, while the baseline errors scatter in sign.

%%%%%%%%%%%%%%%%%%%%%%%%%%%%%%%%%%%%%%%%%%%%%%%%%%%%%%%%%%%%%%%%%%%%%
%%  table: bandshift
%%%%%%%%%%%%%%%%%%%%%%%%%%%%%%%%%%%%%%%%%%%%%%%%%%%%%%%%%%%%%%%%%%%%%
\begin{table}[htbp]
\centering
\caption{Position of the dominant 600--850\,\wn{} band before and after promoting phosphorus, fluorine, sulfur and chlorine to a polarized triple-zeta basis.}
\label{tab:si-bandshift}
\small
\begin{tabular}{lrrrrrr}
\toprule
Agent & Measured & Baseline & Error & tight P & Error & Shift \\
\midrule
GF & 752 & 728 & -24 & 745 & -7 & +18 \\
GD & 752 & 747 & -5 & 759 & +7 & +12 \\
A-230 & 708 & 723 & +15 & 721 & +13 & -2 \\
A-232 & 708 & 657 & -51 & 662 & -46 & +5 \\
A-234 & 714 & 768 & +53 & 794 & +80 & +26 \\
HD & 733 & 741 & +8 & 744 & +11 & +3 \\
\bottomrule
\end{tabular}

\begin{flushleft}\footnotesize The shift is upward in 5 of the 6 molecules, by 3 to 26\,\wn{}, and essentially absent ($-$1.6\,\wn{}) in the remaining one, while the baseline errors scatter from $-$51 to $+$53\,\wn{}. A correction of fixed sign can only help the subset whose error is negative and of comparable size. GF satisfies that condition; GD, which differs only in the ester substituent and was already within 5\,\wn{} of the measured band, does not, and the shift carries it 7\,\wn{} past that band. Positions are the strongest computed line in the window, compared with the strongest measured band.\end{flushleft}
\end{table}
%%%%%%%%%%%%%%%%%%%%%%%%%%%%%%%%%%%%%%%%%%%%%%%%%%%%%%%%%%%%%%%%%%%%%
\clearpage

\subsection{Sensitivity to the band-pairing rule}
\label{si:pairrule}

The frequency scale is fitted by pairing measured band maxima with scaled harmonic
wavenumbers, and this work pairs them by one-to-one assignment, for the reason given in
Section~\ref{si:bands}: under nearest-neighbor matching a single computed band may be
claimed by several measured bands, so residuals of opposite sign cancel and the pair count
is inflated. Because the refitted scale enters every quantity in this section, we report
what the section looks like under both rules (Table~\ref{tab:si-pairrule}).

The two rules agree on the direction and on the ordering of the large effects, and they
agree that cyclosarin is repaired by both the larger phosphorus basis and the
range-separated functional. They differ on soman. Under one-to-one assignment the larger
phosphorus basis degrades soman's window by 0.161 and leaves its assignment correct; under
nearest-neighbor matching it degrades the window by 0.247 and the assignment flips to a
vesicant. The claim that the correction which repairs cyclosarin also breaks its nearest
structural neighbor holds only under the looser rule, and we therefore do not make it. The
weaker statement the stricter rule supports, that the correction which repairs one agent
degrades the same window most for the two agents nearest it, is the one reported in the
main text.

%%%%%%%%%%%%%%%%%%%%%%%%%%%%%%%%%%%%%%%%%%%%%%%%%%%%%%%%%%%%%%%%%%%%%
%%  table: pairing rule sensitivity
%%%%%%%%%%%%%%%%%%%%%%%%%%%%%%%%%%%%%%%%%%%%%%%%%%%%%%%%%%%%%%%%%%%%%
\begin{table}[htbp]
\centering
\caption{The ten trials under both band-pairing rules. $\Delta$ is the change in the 600--850\,\wn{} cosine; the assignment columns give the open-set top-1 result against the pre-correction 53-member library before and after the correction, with \checkmark{} for a correct assignment and $\times$ for an incorrect one.}
\label{tab:si-pairrule}
\footnotesize
\resizebox{\textwidth}{!}{%
\begin{tabular}{llrrcc}
\toprule
& & \multicolumn{2}{c}{$\Delta$ cosine, 600--850 \wn} & \multicolumn{2}{c}{Assignment} \\
\cmidrule(lr){3-4}\cmidrule(lr){5-6}
Agent & Correction & One-to-one & Nearest & One-to-one & Nearest \\
\midrule
GF & $\omega$B97X-D3 & \textbf{+0.370} & \textbf{+0.306} & $\times$ $\rightarrow$ \checkmark & $\times$ $\rightarrow$ \checkmark \\
GF & tight P/F/S/Cl & \textbf{+0.284} & \textbf{+0.294} & $\times$ $\rightarrow$ \checkmark & $\times$ $\rightarrow$ \checkmark \\
A-232 & $\omega$B97X-D3 & \textbf{+0.116} & +0.046 & $\times$ $\rightarrow$ $\times$ & $\times$ $\rightarrow$ $\times$ \\
A-230 & tight P/F/S/Cl & +0.068 & +0.023 & \checkmark $\rightarrow$ \checkmark & \checkmark $\rightarrow$ \checkmark \\
A-232 & CPCM & +0.040 & +0.038 & $\times$ $\rightarrow$ $\times$ & $\times$ $\rightarrow$ $\times$ \\
A-232 & tight P/F/S/Cl & +0.027 & +0.029 & $\times$ $\rightarrow$ $\times$ & $\times$ $\rightarrow$ $\times$ \\
GF & CPCM & +0.005 & -0.013 & $\times$ $\rightarrow$ $\times$ & $\times$ $\rightarrow$ $\times$ \\
HD & tight P/F/S/Cl & -0.010 & -0.016 & \checkmark $\rightarrow$ \checkmark & \checkmark $\rightarrow$ \checkmark \\
GD & tight P/F/S/Cl & \textbf{-0.161} & \textbf{-0.247} & \checkmark $\rightarrow$ \checkmark & \checkmark $\rightarrow$ \textbf{$\times$} \\
A-234 & tight P/F/S/Cl & \textbf{-0.188} & \textbf{-0.150} & $\times$ $\rightarrow$ $\times$ & $\times$ $\rightarrow$ $\times$ \\
\bottomrule
\end{tabular}
}

\begin{flushleft}\footnotesize The baseline calibration itself differs between the rules, shown here on the pre-correction library so that the two rules are compared on identical inputs: one-to-one assignment gives $s_{\mathrm{low}}$ = 1.018, $s_{\mathrm{high}}$ = 0.988 with a position RMSE of 10.6\,\wn{} over 72 pairs, nearest-neighbor matching gives 1.020 and 0.988 with 12.7\,\wn{} over 81 pairs. The one-to-one values are the ones used throughout this work, including in the identification results of Section~\ref{si:ident}. The single qualitative difference between the rules is soman's assignment under the larger phosphorus basis.\end{flushleft}
\end{table}
%%%%%%%%%%%%%%%%%%%%%%%%%%%%%%%%%%%%%%%%%%%%%%%%%%%%%%%%%%%%%%%%%%%%%
\clearpage

%%%%%%%%%%%%%%%%%%%%%%%%%%%%%%%%%%%%%%%%%%%%%%%%%%%%%%%%%%%%%%%%%%%%%
\section{Geometry correction of GF and A-232}
\label{sec:si-geom}

A GFN2 metadynamic conformer search (CREST, 6\,kcal\,mol$^{-1}$ window, eight threads) was
run for each of the eleven agents from its published B3LYP geometry. Nine agents lie within
1.74\,kcal\,mol$^{-1}$ of the search minimum. GF and A-232 do not, and both were
reoptimized from the search minimum at B3LYP/6-311++G(d,p) with \texttt{TightOpt NumFreq}
and the \texttt{\%elprop Polar 1} block, identical to the rest of the library. A third
calculation, the lowest-energy $Z$ conformer of A-232, separates the conformer
contribution from the configurational one. All three terminated normally with no imaginary
frequency.

\begin{figure}[htbp]
\centering
\includegraphics[width=\textwidth]{step4_results.png}
\caption{(a) Effect of correcting a verified geometry error, resolved by window, with the
calibration held at its pre-correction values. A-232 improves in the 600--850\,\wn{}
window and GF deteriorates in it while gaining as much again in 850--1150\,\wn{};
correcting a structural error does not move the agreement in a consistent direction.
(b) The 600--850\,\wn{} window cosine against heavy-atom count. VX and A-242 have
sixteen heavy atoms each and differ by 0.48 in that window, which removes molecular size
as an explanation: the correlation falls from $r$ = $-$0.71 over the nine agents
available before VX was computed to $r$ = $-$0.41 over all ten.}
\label{fig:si-step4}
\end{figure}

\begin{table}[htbp]
\centering
\caption{The three A-232 structures at B3LYP/6-311++G(d,p), with energies relative to the
true minimum and the window cosines obtained with the calibration held at its
pre-correction values so that only the structure changes.}
\label{tab:si-a232}
\small
\begin{tabular}{lrrrr}
\toprule
Structure & $\Delta E$ (kcal\,mol$^{-1}$) & $\Delta G_{298}$ & 600--850 & Full-range self-cosine \\
\midrule
Published geometry ($Z$, arbitrary conformer) & +9.47 & +9.84 & 0.199 & 0.606 \\
Lowest $Z$ conformer & +5.32 & +5.77 & 0.146 & 0.581 \\
Global minimum ($E$) & 0.00 & 0.00 & \textbf{0.446} & \textbf{0.661} \\
\midrule
\textbf{Conformer contribution} & $-$4.15 & & \textbf{$-$0.053} & $-$0.025 \\
\textbf{Configuration contribution} & $-$5.32 & & \textbf{$+$0.300} & $+$0.080 \\
\bottomrule
\end{tabular}

\begin{flushleft}\footnotesize The improvement is carried entirely by the $E/Z$
configuration. A GFN2 estimate of the same decomposition reverses the ordering, placing
the configurational term at 1.77 and the conformer term at 5.37\,kcal\,mol$^{-1}$; the
semi-empirical surface underestimates the amidine $E/Z$ preference by roughly
3.5\,kcal\,mol$^{-1}$ and should not be used to rank these contributions.\end{flushleft}
\end{table}

\begin{table}[htbp]
\centering
\caption{Window-resolved effect of the geometry correction, with the calibration held at
its pre-correction values ($s_\mathrm{low}$ = 1.018, $s_\mathrm{high}$ = 0.988) so that
the only change is the computed spectrum of the agent named.}
\label{tab:si-geomwindow}
\small
\begin{tabular}{lrrrrrr}
\toprule
Agent & 250--600 & 600--850 & 850--1150 & 1150--1400 & 1400--1800 & Mean \\
\midrule
GF, published & 0.772 & 0.548 & 0.730 & 0.944 & 0.661 & 0.731 \\
GF, global minimum & 0.602 & 0.332 & 0.919 & 0.917 & 0.653 & 0.685 \\
\textbf{Change} & $-$0.169 & \textbf{$-$0.216} & $+$0.189 & $-$0.027 & $-$0.007 & $-$0.046 \\
\midrule
A-232, published & 0.793 & 0.199 & 0.674 & 0.863 & 0.810 & 0.668 \\
A-232, global minimum & 0.856 & 0.446 & 0.736 & 0.704 & 0.798 & 0.708 \\
\textbf{Change} & $+$0.063 & \textbf{$+$0.247} & $+$0.062 & $-$0.159 & $-$0.013 & $+$0.040 \\
\bottomrule
\end{tabular}

\begin{flushleft}\footnotesize A-232 improves in the window the main text identifies and
GF deteriorates in it, while GF gains as much again in the adjacent 850--1150\,\wn{}
window. Correcting a verified structural error does not move the agreement in a
consistent direction.\end{flushleft}
\end{table}

\begin{table}[htbp]
\centering
\caption{Sensitivity of the corrected library to the calibration. The two columns differ
only in whether the frequency scale is refitted after the corrected spectra enter the
library.}
\label{tab:si-calibsens}
\small
\begin{tabular}{lcc}
\toprule
 & Scale held at (1.018, 0.988) & Scale refitted to (1.016, 0.990) \\
\midrule
Molecular top-1, closed set & 13/30 & 13/30 \\
Agents correct (majority) & 5/9 & 4/9 \\
HN1 assignment & EXP\_HN1 (\checkmark) & EXP\_HN3 ($\times$) \\
\bottomrule
\end{tabular}

\begin{flushleft}\footnotesize A 0.2\,\% change in the scale factors flips HN1, whose
pre-correction margin was 0.022. The refitted values are the ones reported throughout,
because the calibration procedure is defined before the identification is computed and is
not permitted to depend on it.\end{flushleft}
\end{table}

\section{Measurement repeatability}
\label{si:repeat}

%%%%%%%%%%%%%%%%%%%%%%%%%%%%%%%%%%%%%%%%%%%%%%%%%%%%%%%%%%%%%%%%%%%%%
%%  table: sessionsim
%%%%%%%%%%%%%%%%%%%%%%%%%%%%%%%%%%%%%%%%%%%%%%%%%%%%%%%%%%%%%%%%%%%%%
Table~\ref{tab:si-sessionsim} gives the pairwise similarities on which this section rests.

\begin{table}[htbp]
\centering
\caption{Measurement repeatability. Pairwise cosine similarities between genuine acquisitions, partitioned by whether the pair is the same agent, two agents of the same toxicological class, or two classes.}
\label{tab:si-sessionsim}
\small
\begin{tabular}{lrr}
\toprule
Pair type & $n$ & Mean cosine \\
\midrule
Same agent, different acquisition & 37 & 0.967 \\
Same class, different agent & 127 & 0.642 \\
Different class & 332 & 0.469 \\
\bottomrule
\end{tabular}

\begin{flushleft}\footnotesize Computed on the 32 genuine acquisitions only; the derived copies present in the delivered files were excluded, which the previous version of this analysis did not do. Mann--Whitney: same agent against same class $p = 4.0e-18$, same class against different class $p = 3.8e-24$. The spectra are determined by the material rather than by the session. All acquisitions were taken within a single hour on one instrument, so day-to-day reproducibility remains untested.\end{flushleft}
\end{table}
%%%%%%%%%%%%%%%%%%%%%%%%%%%%%%%%%%%%%%%%%%%%%%%%%%%%%%%%%%%%%%%%%%%%%
\clearpage

%%%%%%%%%%%%%%%%%%%%%%%%%%%%%%%%%%%%%%%%%%%%%%%%%%%%%%%%%%%%%%%%%%%%%
\section{Code and data availability}

The preprocessing module, the analysis scripts that generate every table and figure in
this work, the calibrated computed library, and the ORCA inputs and outputs are deposited
at [repository URL] and archived at [Zenodo DOI].

Raw acquisition files of the authentic agents are not included in that deposit. They are
analytical data on Schedule 1 chemicals generated in an OPCW designated laboratory, and
the appropriate route for them is submission to the OPCW Central Analytical Database
rather than an open archive; they are available to States Parties and designated
laboratories through that route, and to reviewers on request through the editorial
office. Everything needed to reproduce every number in this work from the deposited
preprocessed spectra is in the open deposit.


\end{document}
